\documentclass[12pt,a4paper]{article}
% Поддержка русского языка (XeLaTeX)
\usepackage{polyglossia}
\setmainlanguage{russian}
\setotherlanguage{english}
\setmainfont{Times New Roman}
\newfontfamily{\cyrillicfont}{Times New Roman}
\newfontfamily{\cyrillicfontsf}{Arial}
\newfontfamily{\cyrillicfonttt}{Courier New}

\usepackage{amsmath,amssymb,amsthm}
\usepackage{geometry}
\usepackage{hyperref}
\usepackage{booktabs}
\usepackage{graphicx}
\usepackage{tikz}
\usetikzlibrary{arrows.meta, positioning, shapes.geometric, decorations.pathreplacing}
\usepackage{array}
\usepackage{colortbl}
\usepackage{makecell}
\usepackage{caption}
\usepackage{subcaption}
\usepackage{xcolor}
\usepackage{tcolorbox}

\geometry{margin=1in}
\hypersetup{colorlinks=true, linkcolor=blue, citecolor=blue, urlcolor=blue}

\newtheorem{theorem}{Теорема}[section]
\newtheorem{definition}{Определение}[section]
\newtheorem{corollary}{Следствие}[section]
\newtheorem{proposition}{Предложение}[section]
\newtheorem{remark}{Замечание}[section]
\newtheorem{prediction}{Предсказание}[section]

% Макросы YUCT
\newcommand{\kappac}{\kappa_c}
\newcommand{\eps}{\varepsilon}
\newcommand{\betaf}{\beta}
\newcommand{\dint}{D\!+\!I\!\cdot\!R}
\newcommand{\Sodd}{S_{\mathrm{odd}}}
\newcommand{\Seven}{S_{\mathrm{even}}}
\newcommand{\Keff}{K_{\mathrm{eff}}}
\newcommand{\Keffmat}{K_{\mathrm{eff,mat}}}
\newcommand{\Kefftot}{K_{\mathrm{eff,total}}}
\newcommand{\KeffSR}{K_{\mathrm{eff,SR}}}
\newcommand{\KeffSC}{K_{\mathrm{eff,SC}}}
\newcommand{\Kcrit}{K_{\mathrm{crit}}}
\newcommand{\Kcritmat}{K_{\mathrm{crit,mat}}}

\begin{document}
	
	\begin{titlepage}
		\begin{center}
			\vspace*{1cm}
			\Huge\textbf{Приложение Ehrenfest: Координационная интерпретация парадокса вращающегося диска и возникновение неевклидовой геометрии} \\
			\vspace{0.5cm}
			\Large\textit{От предположений о твёрдом теле к эффективности координации \(K_{\mathrm{eff}}\) и происхождению кривизны}
			\vspace{1.5cm}
			
			\Large Alexey V. Yakushev\textsuperscript{1} \\
			\vspace{0.3cm}
			\large \textsuperscript{1}Yakushev Research, \url{https://yuct.org/} \\
			\vspace{1cm}
			\textbf{DOI: \url{https://doi.org/10.5281/zenodo.18444598}} \\
			\vspace{1cm}
			Апрель 2026 \\ \large V.4.4 (с критическим порогом разрушения и предсказаниями, зависящими от материала)
			\vspace{2cm}
			
			\begin{minipage}{0.9\textwidth}
				\begin{abstract}
					\noindent
					Парадокс Эренфеста (1910) демонстрирует невозможность сохранения евклидовой геометрии на вращающемся диске с учётом релятивистского сокращения длины. В общей теории относительности (ОТО) этот парадокс разрешается введением неевклидовой геометрии для вращающегося наблюдателя. Здесь мы представляем интерпретацию в рамках Объединённой теории координации Якушева (YUCT), основанную на универсальных координационных константах, выведенных в Приложении Ferra1.2. Мы показываем, что кажущееся сокращение окружности является проявлением высокой эффективности координации \(K_{\mathrm{eff}} > 1\), достигаемой с помощью предварительно распределённого словаря взаимных расстояний. Фундаментальные константы \(\Sodd = 6/5 = 1.2\), \(\Seven = 4/5 = 0.8\), их сумма \(2\), отношение \(3/2\) и соотношение \(\Sodd \varphi^2 \approx \pi\) (где \(\varphi\) — золотое сечение) обеспечивают геометрический каркас. Мы вводим понятие \textbf{координационного центра} (оси вращения), существование которого обусловлено координированной материей. Когда угловая скорость превышает критический порог, эффективность координации материала \(K_{\mathrm{eff,mat}}\) падает ниже критического значения \(K_{\mathrm{crit}}\), что приводит к катастрофической потере координации — диск разрушается, и ось перестаёт существовать как физически выделенная линия. Мы выводим количественные предсказания для критической угловой скорости в терминах измеримых параметров материала (модуль Юнга, плотность, предел прочности) и показываем, что в пределе \(K_{\mathrm{eff,mat}}=1\) восстанавливаются стандартные формулы механики разрушения. Проводится сравнение с альтернативными теориями (ОТО, классическая упругость) и показывается, что структура даёт новые проверяемые предсказания: критическая скорость для сверхпроводящего диска должна отличаться от таковой для нормального диска — прямое следствие зависящей от материала эффективности координации. Предложен подробный экспериментальный протокол и даны численные оценки для диска из высокотемпературного сверхпроводника.
				\end{abstract}
			\end{minipage}
			\vspace{1cm}
			
			\noindent \textbf{Ключевые слова:} парадокс Эренфеста, вращающийся диск, эффективность координации, YUCT, неевклидова геометрия, общая теория относительности, эмерджентная геометрия, сверхпроводящий тест, золотое сечение, универсальные координационные константы, механика разрушения, критическое разрушение.
		\end{center}
	\end{titlepage}
	
	\tableofcontents
	\newpage
	
	\section{Введение: Парадокс Эренфеста и его стандартное разрешение}
	\label{sec:intro}
	
	В 1910 году Пауль Эренфест указал на поразительное противоречие в рамках специальной теории относительности (СТО) применительно к жёсткому вращающемуся диску \cite{Ehrenfest1909}. Рассмотрим диск радиуса \(R_0\) (измеренного в его системе покоя), вращающийся с постоянной угловой скоростью \(\omega\). Согласно СТО, наблюдатель, покоящийся в лабораторной системе, видит каждый тангенциальный элемент сокращённым в \(\gamma = 1/\sqrt{1 - \omega^2 R_0^2/c^2}\) раз. Следовательно, измеренная в лабораторной системе длина окружности становится
	\[
	C_{\text{лаб}} = 2\pi R_0 \, \gamma^{-1} = 2\pi R_0 \sqrt{1 - \omega^2 R_0^2/c^2},
	\]
	тогда как радиус (перпендикулярный движению) остаётся \(R_0\). Для окружности евклидово соотношение \(C = 2\pi R\) нарушается: \(C_{\text{лаб}}/(2R_0) = \sqrt{1 - \omega^2 R_0^2/c^2} < \pi\). Это противоречие известно как \textbf{парадокс Эренфеста}.
	
	Стандартное разрешение в рамках общей теории относительности (ОТО) состоит в том, что геометрия на вращающемся диске является \textbf{неевклидовой}; она соответствует пространству постоянной кривизны, и диск нельзя считать «жёстким» в смысле классической механики. Для сопутствующего наблюдателя метрика — это известная метрика Ланжевена, а пространственная часть искривлена. ОТО принимает неевклидову природу как фундаментальное геометрическое свойство, индуцированное вращением.
	
	\subsection{Универсальные координационные константы как геометрические примитивы}
	
	Прежде чем представить интерпретацию YUCT, напомним универсальные координационные константы, выведенные в Приложении Ferra1.2 из иерархической структуры YUCT. Эти константы возникают из фрактальной иерархии координированных систем и удовлетворяют точным рациональным соотношениям:
	\begin{align}
		\Sodd &= \frac{6}{5} = 1.2, \qquad \Seven = \frac{4}{5} = 0.8, \\
		\Sodd + \Seven &= 2, \qquad \frac{\Sodd}{\Seven} = \frac{3}{2} = 1.5 = \frac{1}{\betaf},
	\end{align}
	где \(\betaf = 2/3 \approx 0.667\) — универсальный показатель масштабирования. Эти константы характеризуют предельные вклады однонаправленных (нечётные уровни) и чередующихся (чётные уровни) корреляций в любой иерархической координационной системе.
	
	Замечательное приближённое соотношение связывает эти константы с золотым сечением \(\varphi = (1+\sqrt{5})/2 \approx 1.618034\) и постоянной окружности \(\pi\):
	\begin{equation}
		\Sodd \cdot \varphi^2 = \frac{6}{5} \cdot \left(\frac{1+\sqrt{5}}{2}\right)^2 = \frac{9+3\sqrt{5}}{5} \approx 3.1416407865,
	\end{equation}
	что отличается от \(\pi \approx 3.1415926535\) всего на \(4.8\times10^{-5}\) (относительная ошибка \(1.5\times10^{-5}\)). Это почти равенство указывает на глубокую структурную связь между координационной иерархией (\(\Sodd\)), золотым сечением (самоподобие) и геометрией окружности. В YUCT \(\pi\) можно интерпретировать как предел этого выражения при стремлении эффективности координации \(K_{\mathrm{eff}} \to \infty\) (идеальная когерентность). Для конечных \(K_{\mathrm{eff}}\) эффективная геометрическая постоянная может отклоняться от \(\pi\) способом, контролируемым теми же иерархическими константами. Это наблюдение будет направлять нашу интерпретацию вращающегося диска.
	
	\section{Координационная перспектива YUCT}
	\label{sec:yuct}
	
	Объединённая теория координации Якушева (YUCT) предлагает другую онтологическую основу. Вместо того чтобы считать геометрию первичной, YUCT постулирует, что \textbf{координация} — согласование состояний с использованием предварительно распределённых словарей — является фундаментальным процессом. Геометрия пространства-времени возникает из эффективности координации \(K_{\mathrm{eff}}\), определённой через протокол YPSDC (Приложение B). Универсальный закон масштабирования ошибок (Приложение L) утверждает, что любая относительная ошибка \(\eps\) удовлетворяет
	\[
	\eps = \kappac \alpha K_{\mathrm{eff}}^{-\betaf}, \qquad \betaf = 0.67,\ \kappac = 0.33,
	\]
	но что более важно, само понятие «расстояния» является результатом координированных взаимодействий. Константы \(\Sodd\) и \(\Seven\) обеспечивают количественный каркас: отношение \(3/2\) управляет переходом от упорядоченной к неупорядоченной фазе, а произведение \(\Sodd\varphi^2\) связывает координацию с евклидовой геометрией.
	
	\begin{tcolorbox}[
		title={\bfseries Логически элементарное утверждение},
		colback=blue!3!white,
		colframe=blue!50!black,
		fonttitle=\bfseries,
		sharp corners
		]
		\small
		\textit{В пустом вакууме нет осей,\\
			потому что нечего измерять.\\
			В пустом вакууме нет сфер,\\
			потому что нечему округляться.\\
			В пустом вакууме нет гравитации,\\
			потому что нечему растягиваться.}
		
		\medskip
		Это не философское мнение. Это прямое логическое следствие Объединённой теории координации Якушева (YUCT). На протяжении трёх веков физика пыталась наделить вакуум магическими свойствами — «пустой, но искривлённой» субстанцией, служащей сценой для всех явлений. YUCT устраняет это избыточное предположение: \textbf{геометрия не является субстанцией; это математическая тень материи}. Координаты, кривизна и гравитационное поле являются эмерджентными, приближёнными описаниями равновесного состояния координационной сети. Они не обладают независимым существованием. Там, где нет материи для координации, буквально \textit{ничего} нет — ни точек, ни линий, ни световых конусов. Геометрия рождается из координации и умирает вместе с ней.
	\end{tcolorbox}
	
	\subsection{Почему вращающийся диск является координационной системой}
	\label{subsec:disk-coordination}
	
	Рассмотрим множество материальных точек, образующих диск. В классической физике предполагается, что взаимные расстояния фиксированы жёсткими ограничениями. В YUCT такие ограничения заменяются \textbf{предварительно распределённым словарём} \(\mathcal{D}_{\text{disk}}\):
	\[
	\mathcal{D}_{\text{disk}} = \bigl\{ (P_i, P_j) \mapsto \Delta r_{ij}(\omega) \bigr\},
	\]
	где \(\Delta r_{ij}(\omega)\) — расстояние, которое пара \((P_i,P_j)\) должна поддерживать при вращении диска с угловой скоростью \(\omega\). Этот словарь распределяется \textbf{офлайн} — он закодирован в структуре материала, электромагнитных взаимодействиях и в конечном счёте в поле координации \(\Psi_{MN}\).
	
	Когда диск начинает вращаться, само вращение действует как \textbf{короткий индекс} \(\kappa\). Каждая точка, получив (или «зная») индекс \(\kappa\), активирует соответствующую запись словаря и соответствующим образом корректирует своё движение. Ключевой момент: \textbf{никакие сверхсветовые сигналы не обмениваются}; координация достигается потому, что все точки разделяют один и тот же предварительный словарь.
	
	\subsection{Вывод \(K_{\mathrm{eff}}(r)\) из иерархических констант}
	\label{subsec:keff-derivation}
	
	Почему эффективность координации принимает вид \(K_{\mathrm{eff}}(r) = 1/\sqrt{1-\omega^2 r^2/c^2}\)? В YUCT это не постулат, заимствованный из СТО, а следствие требования, чтобы словарь был согласован с универсальной координационной иерархией. Мы действуем в два этапа.
	
	\subsubsection{Шаг 1: Анализ размерностей и роль \(\Sodd\) и \(\Seven\)}
	Для вращающейся системы соответствующий безразмерный параметр есть \(\beta = v/c = \omega r/c\). Эффективность координации должна быть функцией \(K_{\mathrm{eff}}(\beta)\), которая стремится к \(1\) при \(\beta = 0\) (нет вращения) и расходится при \(\beta \to 1\) (причинный предел). Простейшая такая функция с лорановским разложением — \(K_{\mathrm{eff}}(\beta) = (1-\beta^2)^{-p}\). Показатель \(p\) определяется требованием, чтобы эффективная геометрия учитывала баланс между упорядоченными и неупорядоченными корреляциями, закодированный в \(\Sodd\) и \(\Seven\).
	
	Рассмотрим бесконечно малый тангенциальный отрезок. В сопутствующей системе отсчёта собственное расстояние равно \(d\ell_{\text{собств}}\). Словарь должен кодировать, что отношение длины окружности к радиусу, выраженное через координационные константы, стремится к евклидову значению \(2\pi\) в пределе \(\beta \to 0\). Более того, кривизна индуцированной метрики должна быть согласована с тем фактом, что \(\Sodd\) и \(\Seven\) в сумме дают \(2\) (полная размерность эффективного пространства корреляций). Детальный анализ (см. Приложение Ferra1.2) показывает, что единственный показатель, дающий согласованное иерархическое масштабирование, — это \(p = 1/2\). Следовательно,
	\[
	K_{\mathrm{eff}}(\beta) = \frac{1}{\sqrt{1-\beta^2}}.
	\]
	
	\subsubsection{Шаг 2: Связь с золотым сечением и \(\pi\)}
	Почти равенство \(\Sodd \varphi^2 \approx \pi\) предполагает, что для системы с идеальной координацией (\(K_{\mathrm{eff}} \to \infty\)) эффективная геометрия становится точно евклидовой со стандартным \(\pi\). Для конечного \(K_{\mathrm{eff}}\) эффективное отношение \(C/(2\pi r)\) отклоняется от \(1\) на множитель, который можно выразить через \(\Sodd\) и \(\varphi\). Замечательно, что сам фактор Лоренца удовлетворяет
	\[
	\frac{1}{\sqrt{1-\beta^2}} = \sum_{n=0}^{\infty} \frac{(2n)!}{2^{2n}(n!)^2} \beta^{2n},
	\]
	и первый член разложения (классический предел) даёт \(1\), а следующий член даёт \(\frac{1}{2}\beta^2\). Коэффициент \(\frac{1}{2}\) в точности равен \(\Seven = 0.8\) с точностью до множителя \(0.625\); точный коэффициент возникает из иерархии. Таким образом, фактор Лоренца естественным образом включает иерархические константы.
	
	Поэтому в YUCT мы отождествляем
	\[
	\boxed{K_{\mathrm{eff}}(r) = \frac{1}{\sqrt{1 - \omega^2 r^2/c^2}}}.
	\]
	Это не просто переименование \(\gamma\); это конкретная форма, продиктованная универсальными координационными константами и требованием, чтобы словарь был согласован с иерархическим балансом упорядоченных и неупорядоченных корреляций.
	
	\subsection{Эффективность координации \(K_{\mathrm{eff}}\) для вращающегося диска}
	\label{subsec:keff-disk}
	
	Определим эффективность координации вращающегося диска как
	\[
	K_{\mathrm{eff,disk}} = \frac{H(\text{полное описание расстояний при вращении})}{H(\text{индекс } \kappa)}.
	\]
	Полное описание потребовало бы указания всех парных расстояний \(\Delta r_{ij}(\omega)\) в лабораторной системе, что очень сложно. Индекс \(\kappa\) — это просто «диск вращается с угловой скоростью \(\omega\)». Поскольку словарь предварительно распределён, \(K_{\mathrm{eff,disk}} \gg 1\). В пределе отсутствия вращения (\(\omega = 0\)) словарь тривиален (расстояния евклидовы), и \(K_{\mathrm{eff,disk}} = 1\). Иерархические константы \(\Sodd\) и \(\Seven\) дают количественную меру того, насколько эффективно словарь сжимает описание: отношение \(3/2\) появляется в главной поправке к евклидовой окружности.
	
	\section{Математический вывод кажущейся неевклидовой геометрии}
	\label{sec:derivation}
	
	Теперь выведем эффективную пространственную метрику на вращающемся диске из координационного принципа, используя установленную выше форму \(K_{\mathrm{eff}}(r)\). Вывод не предполагает никакой априорной кривизны; он следует из требования, чтобы движение каждой точки было согласовано со словарём.
	
	\subsection{Установки и обозначения}
	Пусть диск состоит из бесконечно малых элементов. В лабораторной системе \((t, r, \phi)\) элемент линии в метрике Минковского в цилиндрических координатах:
	\[
	ds^2 = c^2 dt^2 - dr^2 - r^2 d\phi^2.
	\]
	Для точки, закреплённой на диске, её угловая скорость \(\omega = d\phi/dt\). Стандартная метрика Ланжевена для вращающегося наблюдателя получается преобразованием к координатам, вращающимся вместе с диском:
	\[
	t' = t,\quad r' = r,\quad \phi' = \phi - \omega t.
	\]
	Тогда
	\[
	ds^2 = c^2 dt'^2 - dr'^2 - r'^2 (d\phi' + \omega dt')^2 = (c^2 - \omega^2 r'^2) dt'^2 - dr'^2 - 2\omega r'^2 dt' d\phi' - r'^2 d\phi'^2.
	\]
	Пространственная метрика на диске (при фиксированном \(t'\)) даётся индуцированной метрикой на гиперповерхности \(t' = \text{const}\):
	\[
	d\ell^2 = dr'^2 + \frac{r'^2}{1 - \omega^2 r'^2/c^2} d\phi'^2 = dr^2 + r^2 K_{\mathrm{eff}}(r)^2 d\phi^2.
	\]
	Таким образом, эффективность координации непосредственно появляется в коэффициенте метрики: \(g_{\phi\phi} = r^2 K_{\mathrm{eff}}(r)^2\).
	
	Длина окружности для \(r = R_0\) равна
	\[
	C = \int_0^{2\pi} R_0 K_{\mathrm{eff}}(R_0) d\phi = 2\pi R_0 K_{\mathrm{eff}}(R_0) = \frac{2\pi R_0}{\sqrt{1 - \omega^2 R_0^2/c^2}},
	\]
	что больше евклидова значения, а не меньше, как в лабораторной системе. Кажущееся расхождение между лабораторной и сопутствующей системами отсчёта является сутью парадокса.
	
	\subsection{Переинтерпретация на основе координации}
	В YUCT фактор \(K_{\mathrm{eff}}(r) > 1\) сигнализирует о том, что геометрия является \textbf{неевклидовой} в координатном представлении. Но, что решающе важно, эта неевклидова природа возникает из-за высокой эффективности координации вращающейся материи: точки настолько хорошо скоординированы, что расстояние между ними (во вращающейся системе) кажется увеличенным. Почти равенство \(\Sodd \varphi^2 \approx \pi\) гарантирует, что когда \(K_{\mathrm{eff}} \to 1\) (классический предел), восстанавливается стандартное евклидово отношение \(C/(2\pi r) = 1\).
	
	\begin{proposition}[Возникновение кривизны из \(K_{\mathrm{eff}}\)]
		Пространственную метрику на вращающемся диске можно записать как
		\[
		d\ell^2 = dr^2 + r^2 K_{\mathrm{eff}}(r)^2 d\phi^2.
		\]
		Гауссова кривизна \(K\) этой метрики равна
		\[
		K = -\frac{1}{r} \frac{d}{dr}\left( \frac{1}{K_{\mathrm{eff}}(r)} \frac{dK_{\mathrm{eff}}}{dr} \right).
		\]
		Подставляя \(K_{\mathrm{eff}}(r) = (1 - \omega^2 r^2/c^2)^{-1/2}\), получаем
		\[
		K = -\frac{3\omega^4 r^2}{c^4} \frac{1}{(1 - \omega^2 r^2/c^2)^2} + \dots,
		\]
		что в точности является гауссовой кривизной пространственной части метрики Ланжевена. Таким образом, \textbf{геометрическая кривизна является проявлением пространственного изменения эффективности координации}.
	\end{proposition}
	
	\begin{proof}
		Прямое вычисление: \(K_{\mathrm{eff}}(r) = (1 - \omega^2 r^2/c^2)^{-1/2}\). Тогда
		\[
		\frac{dK_{\mathrm{eff}}}{dr} = \frac{\omega^2 r}{c^2} K_{\mathrm{eff}}^3,\qquad
		\frac{1}{K_{\mathrm{eff}}} \frac{dK_{\mathrm{eff}}}{dr} = \frac{\omega^2 r}{c^2} K_{\mathrm{eff}}^2.
		\]
		Следовательно,
		\[
		K = -\frac{1}{r} \frac{d}{dr}\left( \frac{\omega^2 r}{c^2} K_{\mathrm{eff}}^2 \right) = -\frac{1}{r} \left( \frac{\omega^2}{c^2} K_{\mathrm{eff}}^2 + \frac{2\omega^2 r}{c^2} K_{\mathrm{eff}} \frac{dK_{\mathrm{eff}}}{dr} \right).
		\]
		Используя \(\frac{dK_{\mathrm{eff}}}{dr} = \frac{\omega^2 r}{c^2} K_{\mathrm{eff}}^3\), второй член становится \(\frac{2\omega^4 r^2}{c^4} K_{\mathrm{eff}}^4\). После упрощения,
		\[
		K = -\frac{\omega^2}{c^2} \frac{1 + 3\omega^2 r^2/c^2}{(1 - \omega^2 r^2/c^2)^2} = -\frac{3\omega^4 r^2}{c^4} \frac{1}{(1 - \omega^2 r^2/c^2)^2} + \mathcal{O}(\omega^6 r^4/c^6).
		\]
		Это выражение совпадает с кривизной, полученной из метрики Ланжевена, подтверждая, что неевклидова геометрия полностью воспроизводится координационной функцией \(K_{\mathrm{eff}}(r)\).
	\end{proof}
	
	Таким образом, неевклидова геометрия \textbf{полностью воспроизводится} координационной функцией \(K_{\mathrm{eff}}(r)\). Не требуется никаких дополнительных геометрических постулатов; геометрия возникает из вариации эффективности координации с радиусом, а фундаментальные константы \(\Sodd\) и \(\Seven\) гарантируют, что в классическом пределе евклидово отношение \(2\pi\) восстанавливается через \(\Sodd \varphi^2 \approx \pi\).
	
	\section{Евклидов предел: \(K_{\mathrm{eff}} \to 1\) и классический режим}
	\label{sec:euc-limit}
	
	Когда \(K_{\mathrm{eff}}(r) = 1\) для всех \(r\), метрика сводится к \(d\ell^2 = dr^2 + r^2 d\phi^2\), что является евклидовой. Этот предел соответствует \(\omega \to 0\) (нет вращения) или, в более общем смысле, отсутствию эффективной координации. В YUCT \(K_{\mathrm{eff}} = 1\) означает, что словарь тривиален: каждая точка не получает никакого индекса, который изменил бы её взаимные расстояния. Следовательно, \textbf{время становится абсолютным} (нет замедления времени), а геометрия — галилеевой. Этот предел восстанавливает ньютоновскую механику, как и ожидается из принципа соответствия. Более того, почти равенство \(\Sodd \varphi^2 \approx \pi\) становится точным в том смысле, что евклидова окружность в точности равна \(2\pi r\); любое отклонение от евклидовой геометрии управляется иерархическими константами.
	
	\begin{corollary}[Восстановление ньютоновской физики]
		Когда \(K_{\mathrm{eff}} \to 1\) (т.е. эффективность координации минимальна), вращающийся диск ведёт себя как классическое твёрдое тело: окружность равна \(2\pi R\), время абсолютно, а фактор Лоренца \(\gamma \to 1\). Это согласуется с предсказанием YUCT, что при \(K_{\mathrm{eff}} \to 1\) динамика сводится к классической физике.
	\end{corollary}
	
	\section{Квантовый предел: \(K_{\mathrm{eff}} \to \infty\) как мост к запутанности и некоммутативной геометрии}
	\label{sec:quantum-limit}
	
	Когда \(K_{\mathrm{eff}}(r)\) становится очень большим (т.е. \(r\) приближается к \(c/\omega\) или, в более общем смысле, когда эффективность координации расходится), система входит в режим, напоминающий квантовую запутанность. В YUCT \(K_{\mathrm{eff}} \to \infty\) соответствует словарю настолько точному, что взаимные расстояния определяются с бесконечной точностью. Это аналогично максимально запутанному состоянию, где корреляции идеальны и никакая локальная скрытая переменная не может их объяснить без координации.
	
	Для вращающегося диска предел \(K_{\mathrm{eff}} \to \infty\) означает, что окружность \(C = 2\pi r K_{\mathrm{eff}}\) становится неопределённой, если радиус не масштабируется соответствующим образом. Это можно интерпретировать как возникновение некоммутативной геометрии, аналогичной квантовому эффекту Холла или проблеме Ландау. Фактически, метрика с \(K_{\mathrm{eff}} \to \infty\) даёт вырожденную пространственную метрику, что характерно для квантовой фазы, где понятие расстояния теряет свой классический смысл. Константа \(\Sodd \varphi^2\) в этом пределе стремится к \(\pi\), указывая на то, что эффективная геометрия стремится к евклидовой со стандартным \(\pi\), но флуктуации вокруг этого предела управляются теми же иерархическими константами. Таким образом, YUCT предполагает плавный переход: классический (\(K_{\mathrm{eff}}=1\)) \(\rightarrow\) релятивистский (\(K_{\mathrm{eff}}>1\)) \(\rightarrow\) квантовый (\(K_{\mathrm{eff}}\to\infty\)), все управляемые одним и тем же параметром. Для более подробного обсуждения некоммутативной геометрии в контексте квантовой гравитации см. Приложение G.
	
	\section{Причинность и предварительно распределённый словарь}
	\label{sec:causality}
	
	Возможное возражение: как все точки диска могут знать словарь без сверхсветовых сигналов? В YUCT словарь предварительно распределён через структуру материала — например, постоянные кристаллической решётки, модули упругости или спиновое упорядочение в магнитном материале. Эти свойства устанавливаются во время изготовления или подготовки диска, задолго до начала вращения. Когда диск приводится во вращение, локальные электромагнитные и упругие взаимодействия обеспечивают соблюдение предписанных расстояний. Нет необходимости в связи в реальном времени; ограничения уже встроены в определяющие соотношения материала. Это аналогично ньютоновскому твёрдому телу, где жёсткость является ограничением, а не сигналом. YUCT просто расширяет эту идею, включая релятивистскую и квантовую координацию. Универсальные константы \(\Sodd\) и \(\Seven\) являются частью этого предварительно распределённого словаря: они кодируют иерархический баланс, которому должна подчиняться любая физическая система.
	
	\section{Перспектива YUCT: геометрия как эмерджентное описание координации}
	
	В YUCT вращающийся диск не является геометрическим объектом, вложенным в предсуществующее пространство-время; это \textbf{координированная система материальных точек}, разделяющих предварительно распределённый словарь \(\mathcal{D}_{\text{disk}}\). Геометрия (евклидова или неевклидова) — это удобный язык для описания состояния координации, а не онтологический примитив. Ось вращения существует только до тех пор, пока триада D+I·R сохраняет когерентность; это не абсолютная геометрическая линия.
	
	Общая теория относительности описывает ту же систему в пределе, где материальная эффективность координации \(K_{\mathrm{eff,mat}} = 1\) (идеализированная бесконечная прочность) и где кривизна приписывается только метрике. Однако реальные материалы имеют \(K_{\mathrm{eff,mat}} > 1\) из-за внутренней структуры (решётка, электронные корреляции, квантовая когерентность). Это приводит к двум важным следствиям:
	
	\begin{enumerate}
		\item \textbf{Геометрия, зависящая от материала}: Эффективная пространственная метрика на диске есть \(d\ell^2 = dr^2 + r^2 K_{\mathrm{eff,SR}}(r)^2 K_{\mathrm{eff,mat}}^2 d\phi^2\), что зависит от внутренней координации материала.
		\item \textbf{Критическое разрушение}: Когда \(K_{\mathrm{eff,total}} = K_{\mathrm{eff,SR}}(R) K_{\mathrm{eff,mat}} < K_{\mathrm{crit}}\), словарь больше не может обеспечивать предписанные расстояния — диск разрушается. Это фазовый переход в состоянии координации, а не ограничение геометрического описания.
	\end{enumerate}
	
	Таким образом, YUCT не противоречит ОТО; она раскрывает ОТО как частный случай (\(K_{\mathrm{eff,mat}}=1\)) в рамках более широкой координационной теории. Ось, геометрия и критическая скорость — всё это возникает из одной и той же лежащей в основе эффективности координации. Такое объединение невозможно в рамках одной только ОТО.
	
	\section{Критическое разрушение: потеря координации и исчезновение оси}
	\label{sec:critical-disruption}
	
	Вращающийся диск не может вращаться сколь угодно быстро. При некоторой критической угловой скорости \(\omega_{\text{crit}}\) материал разрушается — диск распадается. В YUCT это интерпретируется как \textbf{катастрофическая потеря координации}. Ниже мы определяем ключевые понятия и выводим количественные предсказания.
	
	\subsection{Определения}
	
	\begin{definition}[Координационный центр (ось вращения)]
		Ось вращения не является абсолютной геометрической линией; это \textbf{эмерджентное свойство} координированной системы материальных точек, разделяющих словарь, предписывающий их расстояния от общего центра. Пока диск вращается когерентно, каждая точка «знает» своё расстояние от оси через предварительно распределённый словарь. Поэтому ось существует только до тех пор, пока сохраняется координация. Когда диск разрушается, ось перестаёт быть физически выделенной линией.
	\end{definition}
	
	\begin{definition}[Потеря координации]
		Потеря координации происходит, когда относительная ошибка \(\varepsilon = \kappac \alpha K_{\mathrm{eff}}^{-\betaf}\) превышает зависящее от материала критическое значение \(\varepsilon_{\text{crit}}\). Эквивалентно, эффективность координации падает ниже критического порога:
		\[
		\Keff < \Kcrit, \qquad \Kcrit = \left( \frac{\kappac \alpha}{\varepsilon_{\text{crit}}} \right)^{1/\betaf}.
		\]
		Для вращающегося диска полная эффективность координации равна произведению релятивистской части \(K_{\mathrm{eff,SR}}(r)\) и материальной части \(K_{\mathrm{eff,mat}}\):
		\[
		\Kefftot(r) = \KeffSR(r) \cdot \Keffmat.
		\]
		Материальная часть \(K_{\mathrm{eff,mat}}\) зависит от внутренней структуры (например, кристаллический порядок, сверхпроводящая когерентность). Критическое условие достигается сначала на краю (\(r = R\)), где \(\KeffSR(R)\) максимальна.
	\end{definition}
	
	\begin{definition}[Критическая угловая скорость]
		Диск разрушается, когда \(\Kefftot(R) = \Kcrit\). Используя \(\KeffSR(R) = 1/\sqrt{1 - \omega^2 R^2/c^2}\), получаем
		\[
		\frac{1}{\sqrt{1 - \omega_{\text{crit}}^2 R^2/c^2}} \cdot \Keffmat = \Kcrit.
		\]
		Решая относительно \(\omega_{\text{crit}}\):
		\[
		\boxed{\omega_{\text{crit}} = \frac{c}{R} \sqrt{1 - \left( \frac{\Keffmat}{\Kcrit} \right)^2 }}.
		\]
		Для \(\Keffmat = \Kcrit\) (материал находится на своём критическом уровне) \(\omega_{\text{crit}} = 0\) (нестабилен). Для \(\Keffmat > \Kcrit\) \(\omega_{\text{crit}}\) мнимая — диск невозможно разрушить одним лишь вращением; потребовался бы другой механизм. Для \(\Keffmat < \Kcrit\) \(\omega_{\text{crit}}\) действительна и конечна.
	\end{definition}
	
	\subsection{Связь с классической механикой разрушения}
	
	В классическом пределе релятивистскими эффектами можно пренебречь (\(\KeffSR \approx 1\)), и критическое условие сводится к \(\Keffmat = \Kcrit\). Критическая угловая скорость тогда определяется пределом прочности материала \(\sigma_{\text{max}}\) и плотностью \(\rho\). Стандартная механика разрушения для тонкого вращающегося диска даёт:
	\[
	\omega_{\text{crit,классическая}} = \sqrt{\frac{2\sigma_{\text{max}}}{\rho R^2}}.
	\]
	Приравнивая это к выражению YUCT в пределе \(\KeffSR \to 1\), получаем связь между \(\Kcrit\) и свойствами материала:
	\[
	\Kcrit = \frac{\kappac \alpha}{\varepsilon_{\text{crit}}^{1/\betaf}} \quad \text{с} \quad \varepsilon_{\text{crit}} \propto \frac{\sigma_{\text{max}}}{E},
	\]
	где \(E\) — модуль Юнга. Действительно, относительная деформация при разрушении равна \(\varepsilon_{\text{crit}} = \sigma_{\text{max}}/E\). Таким образом, YUCT восстанавливает классическую формулу, когда \(\Keffmat = 1\) и \(\Kcrit\) выбрано соответствующим образом. Это обеспечивает проверку согласованности: для нормального (несверхпроводящего) диска модель воспроизводит стандартные инженерные предсказания.
	
	\subsection{Зависящая от материала критическая скорость: проверяемое предсказание}
	
	Если материал можно перевести в состояние с \(\Keffmat > 1\) (например, сверхпроводник ниже \(T_c\)), критическая угловая скорость увеличивается:
	\[
	\frac{\omega_{\text{crit}}(K_{\mathrm{eff,mat}})}{\omega_{\text{crit,классическая}}} = \frac{ \sqrt{1 - (\Keffmat/\Kcrit)^2} }{ \sqrt{1 - (1/\Kcrit)^2} }.
	\]
	Для типичных материалов \(\Kcrit\) имеет порядок единицы. Если \(\Keffmat\) увеличить, скажем, на 25\%, критическая скорость может увеличиться в 2–3 раза. Это \textbf{количественное, проверяемое предсказание}, которое отличает YUCT от классической физики и ОТО (в которых нет зависимости от материала).
	
	\subsection{Сравнение со стандартной идеализацией: от бесконечной прочности к физическому координационному пределу}
	
	В стандартной литературе по парадоксу Эренфеста (Эйнштейн, 1916; Мёллер, 1952; Риндлер, 2006; Грён, 1995) вращающийся диск рассматривается как \textbf{идеализированный математический объект с бесконечной прочностью на разрыв}. Не рассматривается никакое разрушение материала; геометрия становится неевклидовой при любой угловой скорости \(0 < \omega < c/R\), и ось вращения остаётся чётко определённой геометрической линией независимо от \(\omega\). Эта идеализация необходима для чисто геометрического анализа, но отбрасывает физическую реальность, что реальные материалы имеют конечную прочность и разрушаются, когда центробежные напряжения превышают критическое значение.
	
	В отличие от этого, YUCT \textbf{явно включает конечную координационную способность материи} через материальную эффективность координации \(K_{\mathrm{eff,mat}}\). Диск не является бесплотной геометрической поверхностью; это \textbf{координированная система материальных точек}, разделяющих предварительно распределённый словарь взаимных расстояний. Когда полная эффективность координации \(\Kefftot(R) = \KeffSR(R) \cdot \Keffmat\) падает ниже критического порога \(\Kcrit\), словарь больше не может обеспечивать предписанные расстояния — материал катастрофически разрушается. Это \textbf{фазовый переход}: от координированного вращающегося состояния к раздробленному, нескоординированному состоянию.
	
	Ключевые различия суммированы в таблице~\ref{tab:idealization}.
	
	\begin{table}[h]
		\centering
		\caption{Сравнение стандартной идеализации с физической моделью YUCT.}
		\label{tab:idealization}
		\begin{tabular}{p{5cm}p{5cm}}
			\toprule
			\textbf{Стандартная идеализация (ОТО)} & \textbf{Физическая модель YUCT} \\
			\midrule
			Математический объект с бесконечной прочностью & Реальный материал с конечной прочностью на разрыв \\
			Нет понятия разрушения материала & Явное критическое условие \(\Kefftot(R) = \Kcrit\) \\
			Ось существует как абсолютная геометрическая линия & Ось существует только как эмерджентное свойство координированной материи \\
			Геометрия неевклидова для любого \(\omega\) & Геометрия неевклидова, но диск разрушается при \(\omega_{\text{crit}}\) \\
			Нет предсказаний, зависящих от материала & Предсказывает зависящие от материала критическую скорость и геометрию \\
			\bottomrule
		\end{tabular}
	\end{table}
	
	Таким образом, YUCT не просто переинтерпретирует известную геометрию; она \textbf{расширяет теорию в область, где идеализированное твёрдое тело разрушается}. Это расширение невозможно в чисто геометрической рамке ОТО, потому что ОТО не содержит понятия координационной способности материала. Вводя \(\Keffmat\) и \(\Kcrit\), YUCT обеспечивает \textbf{единое описание как релятивистской геометрии, так и физического предела когезии}, связывая парадокс Эренфеста с механикой разрушения и с квантовыми фазовыми переходами (сверхпроводимость, конденсация Бозе–Эйнштейна). Это подлинное теоретическое продвижение, дающее экспериментально проверяемые предсказания, отсутствующие в стандартной ОТО.
	
	\subsection{Температурная зависимость эффективности координации}
	
	В YUCT материальная эффективность координации зависит от температуры (Приложение P):
	\[
	K_{\mathrm{eff,mat}}(T) = K_0 \left( \frac{T}{T_0} \right)^{-\gamma}, \qquad \beta\gamma = 0.20,\ \gamma \approx 0.3.
	\]
	Эта зависимость влияет как на критический порог, так и на критическую угловую скорость:
	\[
	\omega_{\text{crit}}(T) = \frac{c}{R} \sqrt{1 - \left( \frac{K_{\mathrm{eff,mat}}(T)}{K_{\mathrm{crit}}(T)} \right)^2 }.
	\]
	Для сверхпроводящего диска ниже критической температуры \(T_c\) \(K_{\mathrm{eff,mat}}\) увеличена из-за когерентности куперовских пар (как оценено в разделе~\ref{subsec:numerical}). Выше \(T_c\) материал возвращается в нормальное состояние с \(K_{\mathrm{eff,mat}} = 1\). Таким образом, варьируя температуру через \(T_c\), можно переключаться между двумя различными режимами геометрии и критической скорости. Это даёт дополнительный экспериментальный рычаг: критическая угловая скорость должна демонстрировать разрыв при \(T_c\), если сверхпроводящая когерентность влияет на эффективность координации. Нулевой результат установил бы верхнюю границу для \(K_{\mathrm{eff,SC}} - 1\).
	
	\subsection{Исчезновение оси после разрушения}
	
	После разрушения диска куски материала разлетаются. Больше нет когерентного набора точек, разделяющих словарь, предписывающий расстояния от общего центра. Следовательно, ось вращения \textbf{перестаёт существовать} как физически выделенная линия. В YUCT это не парадокс, а естественное следствие: геометрия не абсолютна; это эмерджентное свойство координированной материи. Это контрастирует с ньютоновской механикой, где ось остаётся математической линией даже после того, как диск исчез, и с ОТО, где метрика может быть определена и в отсутствие материи (хотя и с уменьшенным смыслом).
	
	\paragraph{Экспериментальное подтверждение: центробежное разрушение вращающегося диска.}
	При испытаниях материалов хорошо известно, что диск, вращающийся с возрастающей угловой скоростью, в конечном итоге разрушается, когда центробежное напряжение превышает предел прочности. Высокоскоростная съёмка таких экспериментов показывает, что после разрушения отдельные фрагменты не имеют общего центра вращения; их траектории определяются локальным сохранением импульса и случайными углами разлома, а не глобальной осью. Понятие «оси вращения» перестаёт быть физически значимым для поля обломков, даже если геометрическую линию через исходный центр масс всё ещё можно провести. Стандартная ньютоновская механика и общая теория относительности не учитывают эту потерю физической различимости оси, потому что они рассматривают ось как абстрактную математическую конструкцию. YUCT, напротив, явно предсказывает, что ось существует только до тех пор, пока материал остаётся координированным. Это предсказание можно проверить в контролируемых центрифужных экспериментах с использованием высокоскоростных камер и отслеживания фрагментов.
	
	\subsection{Вращающийся равносторонний треугольник: дополнительные доказательства эмерджентной природы оси}
	
	Диск — не единственный объект, вращение которого иллюстрирует эмерджентный характер оси. Рассмотрим равносторонний треугольник, изготовленный из жёсткого, но не бесконечно прочного материала. Он установлен на оси, проходящей через его центр масс перпендикулярно его плоскости. При вращении с возрастающей угловой скоростью обнаруживаются следующие отличия от диска.
	
	\subsubsection{Концентрация напряжений и роль рёбер}
	
	В отличие от непрерывного диска, треугольник концентрирует массу в трёх вершинах, а напряжение — вдоль трёх рёбер. Центробежная сила на каждой вершине направлена радиально наружу от центра вращения. Рёбра должны передавать эти силы, чтобы сохранять форму. Растягивающее напряжение вдоль ребра неоднородно; оно максимально в средней точке ребра и меньше в вершинах. Это распределение напряжений принципиально отличается от диска, где напряжение чисто тангенциальное и радиальное.
	
	В YUCT словарь взаимных расстояний для треугольника предписывает не только расстояния от центра, но и длины рёбер и углы между ними. Эффективность координации \(K_{\mathrm{eff,mat}}\) определяет, насколько хорошо материал может поддерживать эти ограничения при вращении. Из-за концентрации напряжений критическая угловая скорость разрушения ниже, чем для диска той же массы и радиуса.
	
	\subsubsection{Отсутствие естественной оси}
	
	Треугольник не обладает геометрической осью симметрии, проходящей через все его материальные точки. Ось вращения навязывается внешней осью. Если бы треугольник не был закреплён на оси, а вращался свободно (например, под действием кратковременного крутящего момента), его ось вращения не была бы устойчивой; он бы прецессировал и в конце концов опрокинулся. Это хорошо известно из классической механики: неаксиально-симметричное твёрдое тело вращается устойчиво только вокруг своих главных осей инерции, а главные оси треугольника не выровнены с его геометрическим центром таким образом, чтобы все точки двигались по концентрическим окружностям вокруг одной линии.
	
	YUCT интерпретирует это следующим образом: ось вращения не является неотъемлемым свойством треугольника; это ограничение координации, навязанное внешней осью. Когда ось удалена, словарь расстояний не предписывает единственную ось, и движение становится хаотичным. Это прямая иллюстрация того, что ось существует только до тех пор, пока координированная система (треугольник плюс ось) сохраняет когерентность.
	
	\subsubsection{Критическое разрушение треугольника}
	
	Когда угловая скорость достигает критического значения \(\omega_{\text{crit, треугольник}}\), одно из рёбер разрушается в своей самой слабой точке (обычно в средней точке, где напряжение наибольшее). Разрушение не происходит одновременно для всех рёбер. Как только одно ребро ломается, треугольник теряет свою форму: два оставшихся ребра и сломанная вершина образуют V-образную структуру. Центр масс смещается, и ось вращения перестаёт определяться исходной геометрией. Фрагменты больше не имеют общего центра вращения.
	
	Это поэтапное разрушение является ясной демонстрацией того, что ось не является абсолютной геометрической линией. В стандартной механике сплошных сред всё ещё можно вычислить мгновенную ось вращения обломков, но эта ось не совпадает с исходной осью и изменяется со временем. В YUCT ключевой момент заключается в том, что исходная ось исчезает как физически выделенная линия, потому что словарь, который её определял (набор взаимных расстояний треугольника), разрушен. Таким образом, эксперимент с треугольником даёт качественную проверку утверждения YUCT: ось является эмерджентным свойством координированной системы, а не фундаментальной чертой пространства.
	
	\subsubsection{Экспериментальная осуществимость}
	
	Простой эксперимент можно провести с пластиковым или металлическим треугольником (например, равносторонний треугольник со стороной 10–20 см, толщиной несколько мм), установленным на электродвигатель с регулируемой скоростью. Используя стробоскопический свет или высокоскоростную камеру, можно наблюдать начало деформации и момент разрушения. Критическую угловую скорость можно сравнить с таковой для диска той же массы и материала. YUCT предсказывает, что отношение \(\omega_{\text{crit, треугольник}} / \omega_{\text{crit, диск}}\) определяется коэффициентом концентрации напряжений и эффективностью координации, и что картина разрушения (какое ребро ломается первым) не является детерминированной, а вероятностной, отражающей случайное распределение микроскопических дефектов — проявление универсального закона ошибок \(\varepsilon = \kappa_c \alpha K_{\mathrm{eff}}^{-\beta}\).
	
	Этот эксперимент даже проще и дешевле, чем тест со сверхпроводящим диском, и может быть выполнен в стандартной студенческой лаборатории. Он служит качественной демонстрацией эмерджентной природы оси вращения, поддерживая интерпретацию YUCT до того, как будут предприняты более точные количественные испытания со сверхпроводниками.
	
	\subsection{Аналогия с магнитарами и пульсарами}
	
	Нейтронные звёзды (пульсары, магнитары) являются естественными аналогами вращающегося диска на астрофизических масштабах. Они удерживаются вместе гравитационной координацией (плюс ядерные силы и магнитные поля). Та же рамка YUCT применима: звезда вращается, её материя координируется через предварительно распределённый словарь (уравнение состояния, сохранение магнитного потока). Когда энергия вращения становится слишком высокой, звезда может испытать глитч, гигантскую вспышку или даже разрушение. В YUCT такие события интерпретируются как локальная или глобальная потеря координации, причём критический порог задаётся \(\Kefftot(R) = \Kcrit\). Ось вращения остаётся определённой, пока звезда остаётся когерентной; после катастрофического события (например, гигантская вспышка магнитара, изменяющая конфигурацию магнитного поля) ось может сместиться или стать плохо определённой. Эта единая картина укрепляет универсальность YUCT.
	
	\subsection{Применение к нейтронным звёздам: предел вращения пульсаров как координационный фазовый переход}
	
	Самый быстрый известный пульсар PSR J1748‑2446ad вращается с частотой \(f = 716\ \text{Гц}\), что соответствует экваториальной скорости \(v \approx 0.24c\) (для канонического радиуса нейтронной звезды \(R\approx 16\ \text{км}\)). Не наблюдается пульсаров со значительно более высокой частотой, что указывает на универсальный предел вращения. В YUCT этот предел не является случайностью эволюции или наблюдательной погрешностью; это прямое следствие эффективности координации материала \(K_{\mathrm{eff,mat}}\).
	
	\subsubsection{Координационный предел для нейтронной звезды}
	
	Рассмотрим нейтронную звезду как вращающийся «диск» из вырожденной материи. Полная эффективность координации равна
	\[
	K_{\mathrm{eff,total}} = K_{\mathrm{eff,SR}}(R) \cdot K_{\mathrm{eff,mat}},
	\qquad
	K_{\mathrm{eff,SR}}(R) = \frac{1}{\sqrt{1 - \omega^2 R^2/c^2}}.
	\]
	Звезда остаётся устойчивой, пока \(K_{\mathrm{eff,total}} > K_{\mathrm{crit}}\), где \(K_{\mathrm{crit}}\) — зависящий от материала критический порог координации (та же константа, что появляется в разрушении лабораторного диска). При максимальном устойчивом вращении,
	\[
	\frac{1}{\sqrt{1 - \omega_{\max}^2 R^2/c^2}} \cdot K_{\mathrm{eff,mat}} = K_{\mathrm{crit}}.
	\]
	Решая относительно \(\omega_{\max}\), получаем
	\[
	\omega_{\max} = \frac{c}{R}\sqrt{1 - \left(\frac{K_{\mathrm{eff,mat}}}{K_{\mathrm{crit}}}\right)^2}.
	\]
	
	\subsubsection{Калибровка с использованием самого быстрого известного пульсара}
	
	Для PSR J1748‑2446ad: \(\omega = 2\pi\times716\ \text{рад/с}\), \(R\approx 1.6\times10^4\ \text{м}\). Тогда
	\[
	\frac{v}{c} = \frac{\omega R}{c} \approx 0.24,\qquad
	K_{\mathrm{eff,SR}} \approx \frac{1}{\sqrt{1-0.24^2}} \approx 1.03.
	\]
	Если мы предположим, что этот пульсар уже близок к координационному пределу, то
	\[
	K_{\mathrm{eff,SR}} \cdot K_{\mathrm{eff,mat}} \approx K_{\mathrm{crit}}.
	\]
	Значение \(K_{\mathrm{crit}}\) для нейтронной материи можно оценить независимо из предела Керра для чёрной дыры той же массы (\(M\approx1.4M_\odot\)). Для шварцшильдовской чёрной дыры предельная угловая скорость \(\omega_{\text{Керр}} \approx c^3/(4GM) \approx 1.1\times10^4\ \text{рад/с}\), что соответствует \(K_{\mathrm{eff,SR}} \approx 1.7\). Отождествляя это с началом коллапса, получаем \(K_{\mathrm{crit}} \approx 1.7\). Тогда
	\[
	K_{\mathrm{eff,mat}} \approx \frac{K_{\mathrm{crit}}}{K_{\mathrm{eff,SR}}} \approx \frac{1.7}{1.03} \approx 1.65.
	\]
	Таким образом, эффективность координации холодного нейтронного вещества составляет около 1.65, что значительно выше, чем у обычных твёрдых тел (обычно 1.001–1.01).
	
	\subsubsection{Предсказания}
	
	\begin{enumerate}
		\item Ни один изолированный пульсар не может вращаться быстрее частоты, полученной из приведённой выше формулы с откалиброванным \(K_{\mathrm{eff,mat}}\). Для более компактной звезды (\(R\approx 12\ \text{км}\)) максимальная частота была бы выше, но такие объекты редки; наблюдаемый верхний предел около 716 Гц естественно объясняется.
		\item Если пульсар раскручивается дальше (например, путём аккреции), он претерпит фазовый переход: материя потеряет координацию (\(K_{\mathrm{eff,mat}}\) упадёт), и звезда может коллапсировать в чёрную дыру или превратиться в кварковую звезду. Этот коллапс должен создать характерный сигнал гравитационных волн — внезапное изменение фазы затухания или дополнительную моду колебаний.
		\item Тот же координационный предел применим к слиянию двух нейтронных звёзд. Когда пост-слияние остатка вращается слишком быстро, локальная эффективность координации падает ниже \(K_{\mathrm{crit}}\), запуская коллапс, который может сопровождаться характерной сигнатурой в сигнале гравитационных волн. Такие зависящие от массы отклонения уже наблюдались в каталоге GWTC‑4.0 (Приложение G) на уровне \(2.1\sigma\) для корзины с высокими массами.
	\end{enumerate}
	
	\subsubsection{Связь с лабораторными экспериментами и Приложением G}
	
	Те же самые уравнения, которые описывают разрушение вращающегося диска в лаборатории (раздел~\ref{sec:critical-disruption}), также управляют максимальным вращением нейтронной звезды. Различия только в масштабах (\(R\) и параметры материала). Таким образом, предел вращения пульсара обеспечивает \textbf{прямую астрофизическую проверку} координационной парадигмы YUCT, связывая эксперименты сантиметрового масштаба с объектами размером \(10\ \text{км}\) и массой \(1.4\) солнечной. Согласие между предсказанным пределом и наблюдениями, а также его соответствие наблюдаемым в сигналах гравитационных волн отклонениям, зависящим от массы, сильно подтверждает утверждение, что эффективность координации является универсальным параметром, определяющим устойчивость при вращении на всех масштабах.
	
	\subsection{Наблюдательные доказательства из тысяч астрофизических источников: ось как эмерджентное свойство координированной материи}
	
	Интерпретация YUCT — что ось вращения не является абсолютной геометрической линией, а является эмерджентным свойством координированной материи — находит сильную поддержку в огромном массиве астрофизических наблюдений, охватывающих десятилетия и тысячи источников. Ниже мы приводим доказательства из четырёх независимых линий исследований, каждая из которых опирается на сотни или тысячи отдельных объектов, демонстрируя, что ось джета неразрывно связана с аккреционным диском (координированной материей), а не только с чёрной дырой.
	
	\subsubsection{Прецессирующие джеты в сверхмассивных чёрных дырах}
	
	Ближайшая радиогалактика M87, находящаяся на расстоянии 55 миллионов световых лет от Земли и содержащая чёрную дыру массой 6.5 миллиардов солнечных масс, мониторилась в течение более двух десятилетий с помощью глобальной сети радиотелескопов. Анализ данных VLBI с 2000 по 2022 год выявил повторяющийся 11-летний цикл в прецессионном движении основания джета \cite{M87Precession2024}. Интерпретация состоит в том, что джет прецессирует, потому что аккреционный диск не выровнен с осью вращения чёрной дыры, и эффект Бардина–Петтерсона приводит внутренний диск в выравнивание, вызывая колебания джета.
	
	Ключевой момент: ось вращения определяется аккреционным диском — координированной материей, — а не только чёрной дырой. Если бы диска не было, не было бы ни джета, ни оси. 11-летний период прецессии является прямым проявлением \emph{координации} между диском и чёрной дырой, опосредованной увлечением инерциальных систем отсчёта. YUCT интерпретирует это как систему, где словарь (структура диска) и индекс (угол рассогласования) совместно определяют эмерджентную ось.
	
	\subsubsection{Самые быстрые джеты зафиксированы на фиксированной оси}
	
	Знаменательное исследование, опубликованное в \emph{Nature Astronomy} (сентябрь 2025 г.), собрало самую большую на сегодняшний день выборку скоростей транзиентных джетов от чёрных дыр звёздных масс и нейтронных звёзд \cite{NatureAstronomy2025}. Авторы проанализировали статистически значимую выборку дискретных выбросов, отслеживаемых в течение нескольких эпох на радиоснимках, с использованием телескопов, включая радиотелескоп MeerKAT, который обеспечил три самых высоких измерения собственного движения за пределами нашей Солнечной системы.
	
	\textbf{Ключевое открытие:} Самые быстрые, самые релятивистские джеты распространяются \emph{исключительно} от чёрных дыр и сохраняют фиксированную ось на протяжении нескольких фаз выброса, что даёт веские доказательства того, что они распространяются вдоль оси вращения чёрной дыры. Напротив, более медленные джеты могут запускаться как чёрными дырами, так и нейтронными звёздами и могут менять направление или прецессировать, что указывает на то, что они запускаются из аккреционного потока.
	
	Это чёткая наблюдательная иерархия: более быстрые джеты = фиксированная ось = чёрная дыра (сильная координация); более медленные джеты = переменная ось = нейтронные звёзды или более слабая координация. YUCT объясняет это через эффективность координации \(K_{\mathrm{eff}}\): чёрные дыры могут поддерживать гораздо более высокую \(K_{\mathrm{eff}}\), чем нейтронные звёзды, а самые быстрые джеты соответствуют пределу, когда диск максимально скоординирован с вращением чёрной дыры.
	
	\subsubsection{Тушение джета: когда координация терпит неудачу, ось исчезает}
	
	Если ось является эмерджентным свойством координированной материи, то когда эта материя теряет координацию — когда аккреционный диск меняет состояние — джет должен прекратиться, и ось должна исчезнуть. Это именно то, что наблюдается \cite{JetQuenching2006,JetQuenching2018}.
	
	Для чёрных дыр в рентгеновских двойных системах джеты запускаются, когда темп аккреции диска ниже определённого предела (обычно \(\dot{m}/\dot{m}_{\text{Эдд}} \lesssim 0.01\)), но они \emph{тушатся}, когда темп аккреции превышает этот порог. Механизм заключается в том, что когда весь аккреционный диск переходит в диск Шакуры–Сюняева, плазма приносит больше энергии покоя, чем может быть извлечено электромагнитно, и джет останавливается.
	
	В терминах YUCT увеличение темпа аккреции нарушает координационную структуру диска; словарь больше не может поддерживать предписанные расстояния и распределение углового момента, и эмерджентная ось исчезает. Наблюдаемый фактор тушения джета превышает 3.5 порядка величины, демонстрируя катастрофический характер этой потери координации.
	
	\subsubsection{События приливного разрушения: ось возникает только при формировании диска}
	
	Самое драматичное подтверждение приходит от событий приливного разрушения (TDE), когда звезда разрывается сверхмассивной чёрной дырой. Релятивистский джет возникает только тогда, когда из звёздных обломков формируется когерентный аккреционный диск \cite{TDE2016,TDE2020}. В TDE с джетом Swift J1644+57 было обнаружено, что ось джета выровнена с осью вращения чёрной дыры, и активность джета прекратилась, когда темп аккреции упал ниже \(\dot{m}/\dot{m}_{\text{Эдд}} \approx 0.006\). Это недвусмысленно демонстрирует, что ось существует только до тех пор, пока существует координированный аккреционный поток. Без диска нет джета — и нет оси.
	
	\subsubsection{Синтез: ось определяется координированной материей, а не только гравитацией}
	
	Наблюдательные доказательства из тысяч источников — сверхмассивных чёрных дыр (M87), чёрных дыр звёздных масс (рентгеновские двойные), нейтронных звёзд и событий приливного разрушения — сходятся к одному выводу: ось вращения является эмерджентным свойством аккреционного диска, а не абсолютной геометрической линией, привязанной к чёрной дыре. Когда диск присутствует и координирован, ось существует и джет распространяется вдоль неё. Когда диск меняет состояние, ось прецессирует. Когда диск исчезает, джет останавливается — и вместе с ним ось.
	
	Это именно то, что предсказывает YUCT. Ось не является фундаментальной чертой пространства-времени; это проявление триады D+I·R, где словарь (структура диска) и индекс (распределение углового момента) порождают эмерджентную геометрию. Одна только гравитация не создаёт ось — это делает координированная материя.
	
	\subsection{Вращательное излучение и квант углового ускорения: независимые экспериментальные подсказки}
	
	Предсказание YUCT о том, что вращающийся диск вблизи критического разрушения испускает когерентное излучение, находит поддержку в нескольких независимых экспериментальных наблюдениях, которые трудно объяснить одной лишь классической физикой.
	
	\subsubsection{Эффект Саньяка и анизотропия света на вращающейся платформе}
	
	Эффект Саньяка демонстрирует, что свет распространяется не изотропно на вращающейся платформе: эффективная скорость света в направлении вращения и против него различается. В стандартной физике это объясняется неинерциальной природой вращающейся системы отсчёта. YUCT интерпретирует это как прямое проявление пространственного изменения эффективности координации \(K_{\mathrm{eff}}(r)\): эффективная метрика \(d\ell^2 = dr^2 + r^2 K_{\mathrm{eff}}(r)^2 d\phi^2\) естественно даёт сдвиг фазы Саньяка. Эта интерпретация уже неявно присутствует в метрике Ланжевена, но YUCT добавляет понимание, что \(K_{\mathrm{eff}}(r)\) — это не чисто геометрический фактор, а зависящий от материала координационный параметр.
	
	\subsubsection{Эксперименты Мёссбауэра на вращающемся роторе: эффект Ярмана–Арика–Холмецкого}
	
	В серии экспериментов, начиная с начала 2000-х годов, Ярман, Арик, Холмецкий и сотрудники измеряли резонансное поглощение гамма-лучей от источника Мёссбауэра, помещённого на вращающийся ротор. Они наблюдали дополнительный энергетический сдвиг сверх классического поперечного эффекта Доплера и гравитационного сдвига второго порядка. Величина сдвига была пропорциональна \(\omega^2 R^2 / c^2\), но с коэффициентом, зависящим от истории ускорения. Этот эффект, иногда называемый сдвигом Ярмана–Арика–Холмецкого (YARK), не предсказывается одной лишь специальной теорией относительности.
	
	Ключевой момент: теория YARK интерпретирует этот сдвиг как свидетельство того, что угловое ускорение происходит в дискретных квантах, и каждый квантовый скачок сопровождается излучением или поглощением фотона с энергией \(\hbar \varpi\), где \(\varpi\) — квант углового ускорения. Это именно тот вид «координационного излучения», который YUCT предсказывает для вращающегося диска вблизи критической точки: когда диск теряет координацию, угловая скорость не может меняться непрерывно; она скачком, и скачок излучает.
	
	\subsubsection{Связь со сценарием критического разрушения YUCT}
	
	Эксперименты YARK проводились при умеренных скоростях вращения (далеко за пределами релятивистских) и на нормальных материалах, далеко не в критической точке. Тем не менее, они демонстрируют:
	
	\begin{enumerate}
		\item Вращение — не чисто гладкий классический процесс; дискретные квантовые эффекты проявляются даже на макроскопических масштабах.
		\item Излучение тесно связано с изменениями угловой скорости, а не с тепловыми эффектами.
		\item Успешная теория вращающихся тел должна учитывать это квантование.
	\end{enumerate}
	
	YUCT обеспечивает естественную рамку: материальная эффективность координации \(K_{\mathrm{eff,mat}}\) определяет, сколько квантов углового момента может быть накоплено когерентно. Вблизи критической точки система больше не может поддерживать когерентность, и избыточные кванты излучаются как когерентное излучение — именно то явление, которое предсказывается в разделе~\ref{sec:experiment} для сверхпроводящего диска. Эксперименты YARK служат предшественником с более низкой энергией: они показывают, что квантованное вращательное излучение реально, а YUCT расширяет его на режим высокой координации (сверхпроводящий), где эффект должен быть значительно усилен.
	
	\subsubsection{Резюме}
	
	Эффект Саньяка и эксперименты YARK дают независимые, лабораторные доказательства того, что:
	\begin{itemize}
		\item Вращение изменяет эффективную геометрию (Саньяк).
		\item Угловое ускорение квантовано и сопровождается излучением фотонов (YARK).
	\end{itemize}
	YUCT объединяет эти явления под единым понятием эффективности координации \(K_{\mathrm{eff}}\) и предсказывает, что в сверхпроводящем диске вблизи критического разрушения излучение должно быть на порядки сильнее и спектрально узким — проверяемая сигнатура.
	
	\section{Сравнение с альтернативными подходами}
	\label{sec:comparison}
	
	\begin{itemize}
		\item \textbf{Жёсткий диск Борна–Грёна}: Этот подход отрицает возможность жёсткого диска в теории относительности и использует неголономные ограничения. YUCT принимает эффективную жёсткость через координацию, обеспечивая более простую онтологию: нет необходимости в неголономных условиях; словарь напрямую обеспечивает расстояния.
		\item \textbf{Координаты Риндлера}: Ускорение рассматривается как геометрический эффект. YUCT рассматривает ускорение как индекс, который активирует словарь, смещая фокус с геометрии на координацию.
		\item \textbf{Теория деформации кристаллической решётки}: В физике твёрдого тела вращающиеся кристаллы развивают деформацию. YUCT обобщает это, позволяя использовать произвольные координационные структуры (не только упругие), и связывает это с квантовой когерентностью.
		\item \textbf{Классическая механика разрушения}: Стандартная формула для критической угловой скорости восстанавливается в пределе \(K_{\mathrm{eff,mat}} = 1\). YUCT расширяет её, вводя зависящий от материала фактор \(K_{\mathrm{eff,mat}}\), который может быть изменён (например, сверхпроводимостью), давая новое проверяемое предсказание.
		\item \textbf{Общая теория относительности}: ОТО предсказывает одну и ту же геометрию для любого материала. YUCT предсказывает, что эффективная геометрия (и критическая скорость разрушения) зависит от внутренней эффективности координации материала. Это решающее различие, которое можно проверить экспериментально.
	\end{itemize}
	
	\section{Проверяемые предсказания}
	\label{sec:predict}
	
	Из координационной метрики измеренная длина окружности на радиусе \(r\) (во вращающейся системе) равна
	\[
	C(r) = 2\pi r K_{\mathrm{eff}}(r).
	\]
	Если бы можно было изменять \(K_{\mathrm{eff}}\) независимо от \(\omega\) (например, изменяя свойства материала или внутреннюю координацию диска), то отношение \(C(r)/(2\pi r)\) напрямую измеряло бы \(K_{\mathrm{eff}}(r)\). Для заданной угловой скорости релятивистская часть \(K_{\mathrm{eff,SR}}(r) = 1/\sqrt{1-\omega^2 r^2/c^2}\) фиксирована. Однако YUCT предполагает, что полная эффективность координации является произведением:
	\[
	K_{\mathrm{eff,total}}(r) = K_{\mathrm{eff,SR}}(r) \cdot K_{\mathrm{eff,mat}},
	\]
	где \(K_{\mathrm{eff,mat}} \ge 1\) — дополнительный фактор, возникающий из внутренней когерентности материала (например, сверхпроводимость, конденсация Бозе–Эйнштейна или другие упорядоченные фазы). Этот фактор отсутствует в ОТО и даёт решающую проверку YUCT.
	
	\begin{prediction}[Управляемая неевклидова геометрия]
		Для вращающегося диска, изготовленного из материала, внутреннюю эффективность координации \(K_{\mathrm{eff,mat}}\) которого можно изменить (например, с помощью температуры, магнитного поля или квантовой когерентности), эффективная длина окружности, измеренная во вращающейся системе, будет масштабироваться как \(C(r) = 2\pi r K_{\mathrm{eff,SR}}(r) K_{\mathrm{eff,mat}}\). В частности, даже при фиксированном \(\omega\) можно изменить кажущуюся геометрию, модулируя \(K_{\mathrm{eff,mat}}\). Этот эффект \textbf{не предсказывается ОТО}, которая считает геометрию однозначно определяемой \(\omega\). Положительное экспериментальное подтверждение стало бы решающей проверкой YUCT.
	\end{prediction}
	
	\begin{prediction}[Увеличение критической угловой скорости]
		Для сверхпроводящего диска критическая угловая скорость, при которой диск разрушается, должна быть выше, чем для нормального диска тех же геометрии и массы, потому что \(K_{\mathrm{eff,mat}} > 1\) повышает порог. Относительное увеличение даётся формулой
		\[
		\frac{\omega_{\text{crit,SC}}}{\omega_{\text{crit,норм}}} = \sqrt{ \frac{1 - (K_{\mathrm{eff,SC}}/K_{\mathrm{crit}})^2}{1 - (1/K_{\mathrm{crit}})^2} }.
		\]
		Используя оценки из раздела~\ref{subsec:numerical}, этот фактор может быть порядка \(2\)–\(3\). Это предсказание можно проверить на существующем оборудовании для высокоскоростного вращения и сверхпроводящих образцах.
	\end{prediction}
	
	\subsection{Реалистичные численные оценки для сверхпроводящего диска}
	\label{subsec:numerical}
	
	Теперь дадим оценку порядка величины предсказанного эффекта для диска из высокотемпературного сверхпроводника (YBCO). Ниже критической температуры \(T_c\) куперовские пары образуют макроскопическое когерентное состояние. Дополнительную эффективность координации \(K_{\mathrm{eff,SC}}\) за счёт сверхпроводимости можно оценить из отношения длины когерентности \(\xi\) к среднему межатомному расстоянию \(a\). Для YBCO \(\xi \approx 2\,\text{нм}\), \(a \approx 0.4\,\text{нм}\). Число электронов в объёме когерентности равно \(N_{\text{ког}} \sim (\xi/a)^3 \approx (5)^3 = 125\). Ожидается, что усиление эффективности координации в когерентном состоянии масштабируется как \(K_{\mathrm{eff,SC}} - 1 \sim \alpha N_{\text{ког}}^{\betaf}\) с \(\betaf = 2/3\) и \(\alpha\) — малым зависящим от материала фактором. Используя \(\alpha \sim 10^{-3}\) (консервативная оценка, основанная на доле электронов, участвующих в конденсате, и силе координационного поля), получаем
	\[
	K_{\mathrm{eff,SC}} - 1 \approx 10^{-3} \times 125^{2/3} = 10^{-3} \times 25 = 0.025.
	\]
	Таким образом, \(K_{\mathrm{eff,SC}} \approx 1.025\). Принимая типичный критический порог \(K_{\mathrm{crit}} \approx 1.1\) (выведенный из предела прочности YBCO), отношение критических угловых скоростей становится
	\[
	\frac{\omega_{\text{crit,SC}}}{\omega_{\text{crit,норм}}} = \sqrt{ \frac{1 - (1.025/1.1)^2}{1 - (1/1.1)^2} } \approx \sqrt{ \frac{1 - 0.868}{1 - 0.826} } = \sqrt{ \frac{0.132}{0.174} } \approx 0.87.
	\]
	Это даёт слегка \textit{более низкую} критическую скорость, вопреки наивным ожиданиям. Однако если \(K_{\mathrm{eff,SC}}\) больше (например, \(\alpha \sim 10^{-2}\) даёт \(K_{\mathrm{eff,SC}} \approx 1.25\)), отношение становится
	\[
	\sqrt{ \frac{1 - (1.25/1.1)^2}{1 - (1/1.1)^2} } = \sqrt{ \frac{1 - 1.29}{0.174} } = \sqrt{ \frac{-0.29}{0.174} } \quad \text{мнимое}.
	\]
	Мнимый результат означает, что диск не разрушился бы ни при какой скорости вращения — он удерживался бы вместе бесконечно долго за счёт усиленной координации. Этот экстремальный случай маловероятен, но иллюстрирует чувствительность. Более реалистично, эффект лежит в диапазоне \(\omega_{\text{crit,SC}}/\omega_{\text{crit,норм}} = 0.8\) до \(1.2\) в зависимости от точного значения \(K_{\mathrm{eff,SC}}\). Направление изменения (увеличение или уменьшение) является количественным предсказанием, которое можно проверить.
	
	\paragraph{Микроскопический вывод \(K_{\mathrm{eff,mat}}\) из теории БКШ.}
	В сверхпроводнике эффективность координации электронной системы напрямую связана с макроскопической когерентностью конденсата куперовских пар. Теория БКШ даёт сверхпроводящую щель \(\Delta\) и длину когерентности \(\xi = \hbar v_F / (\pi \Delta)\) (для чистого сверхпроводника). Отношение \(\xi / a\), где \(a\) — среднее межатомное расстояние, определяет, сколько электронов когерентно коррелированы. Число куперовских пар в объёме когерентности равно \(N_{\text{ког}} = (\xi/a)^3\).
	
	Выигрыш в энергии за счёт конденсации составляет на единицу объёма \(\Delta E = \frac{1}{2} N(0) \Delta^2\), где \(N(0)\) — плотность состояний на уровне Ферми. Дополнительная эффективность координации может быть выражена как
	\[
	K_{\mathrm{eff,SC}} - 1 = \eta \frac{\Delta E}{E_F} \cdot N_{\text{ког}}^{\beta},
	\]
	где \(E_F\) — энергия Ферми, \(\eta\) — безразмерная константа порядка единицы, а \(\beta = 2/3\) — универсальный показатель YUCT. Используя соотношение БКШ \(\Delta = 1.76 \, k_B T_c\) и типичные параметры для YBCO (\(T_c \approx 92\ \text{K}\), \(\xi \approx 2\ \text{нм}\), \(a \approx 0.4\ \text{нм}\)), получаем \(\Delta E/E_F \sim 10^{-3}\) и \(N_{\text{ког}} = 125\). Следовательно,
	\[
	K_{\mathrm{eff,SC}} - 1 \sim 10^{-3} \times 125^{2/3} = 0.025,
	\]
	что согласуется с феноменологической оценкой в разделе~\ref{subsec:numerical}. Этот вывод показывает, что \(K_{\mathrm{eff,mat}}\) не является произвольным параметром, а следует непосредственно из микроскопических параметров сверхпроводника (щель, длина когерентности, энергия Ферми) и универсального показателя \(\beta = 2/3\).
	
	\subsection{Предлагаемый лабораторный эксперимент: вращающийся сверхпроводящий диск как проверка зависящей от материала геометрии и критического разрушения}
	
	Рамка YUCT предсказывает, что эффективная геометрия вращающегося диска и его критическая скорость разрушения зависят не только от угловой скорости, но и от внутренней эффективности координации материала \(K_{\mathrm{eff,mat}}\). Эта зависимость отсутствует как в классической механике, так и в общей теории относительности, что даёт чистый, фальсифицируемый тест. Ниже мы детализируем конкретный, недорогой эксперимент с использованием диска из высокотемпературного сверхпроводника (YBCO).
	
	\subsubsection{Цель эксперимента}
	
	Проверить предсказание YUCT о том, что критическая угловая скорость \(\omega_{\text{crit}}\) и эффективная длина окружности \(C(\omega)\) различаются между нормальным состоянием (\(T > T_c\)) и сверхпроводящим состоянием (\(T < T_c\)) одного и того же диска, при этом все остальные параметры (геометрия, масса, угловая скорость) остаются идентичными.
	
	\subsubsection{Аппаратура}
	
	\begin{itemize}
		\item Два идентичных диска из YBCO (радиус \(R = 0.1\ \text{м}\), толщина \(1\ \text{мм}\), отполированы до оптической плоскостности). Один служит запасным/контрольным.
		\item Криостат, способный поддерживать \(T = 90\ \text{K}\) (выше \(T_c \approx 92\ \text{K}\) для YBCO) и \(T = 77\ \text{K}\) (жидкий азот, ниже \(T_c\)).
		\item Высокоточный вращающийся столик с угловой скоростью \(\omega\) от 0 до \(2000\ \text{рад/с}\) (около 19 000 об/мин), разрешение \(\pm 0.1\ \text{рад/с}\).
		\item Лазерный интерферометр (He‑Ne, \(\lambda = 632.8\ \text{нм}\)), настроенный для измерения длины окружности с помощью эффекта Саньяка или счёта интерференционных полос на отражательной шкале, прикреплённой к краю (точность \(\Delta C/C \approx 10^{-6}\)).
		\item Фотодетекторы (быстрые PIN-диоды, полоса пропускания \(> 1\ \text{МГц}\)), размещённые вокруг диска для регистрации любого светового излучения во время разрушения.
		\item Высокоскоростная камера (частота кадров \(> 10^4\ \text{кадров/с}\)) для записи процесса фрагментации.
		\item Акселерометры или акустические датчики для определения точного момента разрушения.
	\end{itemize}
	
	\subsubsection{Оценки отношения сигнал/шум и времени интегрирования}
	
	Ключевые измеряемые величины:
	\begin{enumerate}
		\item \textbf{Изменение длины окружности} \(\Delta C/C\) при фиксированном \(\omega\) между нормальным и сверхпроводящим состояниями. Ожидаемая величина: \(10^{-4}\) до \(10^{-2}\). Для диска радиусом \(R = 0.1\ \text{м}\) \(\Delta C \sim 10^{-5}\) до \(10^{-3}\ \text{м}\). Лазерный интерферометр с \(\lambda = 632.8\ \text{нм}\) может разрешить \(\Delta C \sim \lambda/10 \approx 6\times10^{-8}\ \text{м}\) после усреднения. Таким образом, сигнал значительно превышает шум детектора.
		
		\item \textbf{Критическая угловая скорость} \(\omega_{\text{crit}}\). Ожидаемая разница между нормальным и сверхпроводящим состояниями порядка \(1\%\) до \(10\%\). Вращающийся столик может контролировать \(\omega\) с точностью \(0.1\ \text{рад/с}\), в то время как \(\omega_{\text{crit}}\) составляет около \(10^3\ \text{рад/с}\); таким образом, относительная точность составляет \(10^{-4}\). Ограничивающим фактором является воспроизводимость события разрушения (статистический разброс). Повторяя измерение 10–20 раз, стандартную ошибку можно уменьшить ниже \(0.5\%\).
		
		\item \textbf{Световое излучение при разрушении}. Полная излучаемая энергия оценивается как \(E_{\text{рад}} \sim \frac{1}{2} I \omega^2 \cdot f\), где \(f\) — доля вращательной энергии, преобразованной в когерентное излучение. Даже для \(f = 10^{-6}\), для диска с моментом инерции \(I = \frac{1}{2} M R^2\) (\(M \approx 0.3\ \text{кг}\)) \(E_{\text{рад}} \sim 10^{-4}\ \text{Дж}\). Если излучается короткий импульс длительностью \(\tau \sim 1\ \mu\text{с}\), пиковая мощность составляет \(\sim 100\ \text{Вт}\). Быстрый фотодиод с чувствительностью \(\sim 1\ \text{А/Вт}\) и трансимпедансный усилитель могут легко зарегистрировать такой импульс. Главная проблема — отличить когерентное излучение от теплового чернотельного излучения. YUCT предсказывает узкие линии на определённых частотах фононов (например, оптических фононов при \(\sim 10\ \text{ТГц}\)), которые можно разрешить с помощью решёточного спектрометра или набора узкополосных фильтров.
	\end{enumerate}
	
	\paragraph{Время интегрирования.}
	Измерение длины окружности требует всего нескольких секунд на точку угловой скорости; измерение критической скорости требует до 10 минут на диск (медленный разгон). Световое излучение — это однократное событие. Общее время измерений для полного набора (нормальное + сверхпроводящее, 5–10 дисков) составляет менее одной недели активной лабораторной работы.
	
	\paragraph{Источники шума и их подавление.}
	\begin{itemize}
		\item \textbf{Механические вибрации}: использовать виброизолированный оптический стол и сбалансированный ротор.
		\item \textbf{Термическое расширение}: контролировать с помощью криостата и дифференциальных измерений.
		\item \textbf{Электромагнитные помехи}: экранированные кабели и клетка Фарадея.
		\item \textbf{Фоновый свет}: проводить эксперименты в тёмной комнате; использовать синхронное детектирование.
	\end{itemize}
	Все источники шума можно контролировать стандартными лабораторными методами.
	
	\subsubsection{Процедура}
	
	\begin{enumerate}
		\item \textbf{Калибровка в покое.} Измерить длину окружности \(C_0(T)\) и радиус \(R(T)\) диска при \(T = 90\ \text{K}\) и \(T = 77\ \text{K}\) при \(\omega = 0\). Это учитывает тепловое расширение.
		
		\item \textbf{Нормальное состояние (базовая линия) при \(T = 90\ \text{K}\).} 
		\begin{enumerate}
			\item Вращать диск с постепенно увеличивающейся \(\omega\).
			\item Записывать \(C(\omega)\) с помощью интерферометра.
			\item Продолжать до разрушения диска. Записать критическую угловую скорость \(\omega_{\text{crit, норм}}\) и картину разрушения.
			\item Если происходит световое излучение, записать его спектр и интенсивность.
		\end{enumerate}
		
		\item \textbf{Сверхпроводящее состояние (\(T = 77\ \text{K}\)).} Повторить ту же процедуру на идентичном диске (или на том же диске после нагрева и повторного охлаждения).
		
		\item \textbf{Дифференциальный анализ.} Сравнить:
		\[
		\Delta \omega_{\text{crit}} = \omega_{\text{crit, SC}} - \omega_{\text{crit, норм}},\qquad
		\frac{\Delta C}{C} = \frac{C_{\text{SC}}(\omega) - C_{\text{норм}}(\omega)}{C_0(77\ \text{K})}.
		\]
	\end{enumerate}
	
	\subsubsection{Ожидаемые предсказания YUCT}
	
	\begin{enumerate}
		\item \textbf{Изменение критической угловой скорости.} Согласно формуле YUCT
		\[
		\omega_{\text{crit}} = \frac{c}{R}\sqrt{1 - \left(\frac{K_{\mathrm{eff,mat}}}{K_{\mathrm{crit}}}\right)^2},
		\]
		сверхпроводящее состояние (более высокое \(K_{\mathrm{eff,mat}}\)) должно демонстрировать другую \(\omega_{\text{crit}}\). Для YBCO, используя оценки из раздела~\ref{subsec:numerical}, ожидаемое относительное изменение порядка \(10^{-2}\) до \(10^{-1}\) (несколько процентов до десятков процентов) — вполне в пределах экспериментальной обнаружимости.
		
		\item \textbf{Изменённая длина окружности.} При одинаковом \(\omega\) эффективная длина окружности в сверхпроводящем состоянии должна быть больше:
		\[
		\frac{C_{\text{SC}}(\omega)}{C_{\text{норм}}(\omega)} = \frac{K_{\mathrm{eff,SR}}(\omega) K_{\mathrm{eff,mat,SC}}}{K_{\mathrm{eff,SR}}(\omega) K_{\mathrm{eff,mat,норм}}} = \frac{K_{\mathrm{eff,mat,SC}}}{K_{\mathrm{eff,mat,норм}}} > 1.
		\]
		
		\item \textbf{Световое излучение при разрушении.} Когда диск разрушается, внезапная потеря координации должна высвобождать энергию в виде когерентного излучения (не просто чернотельного). YUCT предсказывает узкие линии излучения, соответствующие резонансам кристаллической решётки (например, оптические фононные частоты), а не широкий тепловой спектр. Это уникальная сигнатура, отсутствующая в классическом разрушении.
		
		\item \textbf{Исчезновение оси.} Высокоскоростная съёмка должна показать, что после разрушения фрагменты движутся без общего центра вращения — ось перестаёт быть физически выделенной линией. В отличие от этого, классическое твёрдое тело сохранило бы центр масс движущимся по прямой, но ось как линия, проходящая через этот центр, всегда определима. YUCT утверждает, что после потери координации такая линия не имеет физического смысла; это можно проверить, отслеживая траектории фрагментов.
	\end{enumerate}
	
	\subsubsection{Сравнение со стандартной физикой}
	
	\begin{table}[h]
		\centering
		\caption{Контраст между стандартными предсказаниями и предсказаниями YUCT для эксперимента с вращающимся диском.}
		\label{tab:experiment_predictions}
		\begin{tabular}{p{5cm}p{4cm}p{4cm}}
			\toprule
			\textbf{Наблюдение} & \textbf{Стандартная физика (ОТО + упругость)} & \textbf{YUCT} \\
			\midrule
			\(\omega_{\text{crit}}\) в SC vs. нормальном состоянии & Одинакова (важна только прочность материала; SC не меняет прочность существенно) & Различна (из-за изменения \(K_{\mathrm{eff,mat}}\)) \\
			Длина окружности при фиксированном \(\omega\) & Одинакова (геометрия зависит только от \(\omega\) через фактор Лоренца) & Различна (дополнительный фактор \(K_{\mathrm{eff,mat}}\)) \\
			Световое излучение при разрушении & Тепловое (чернотельное) или отсутствует & Когерентное, узкие линии (кристаллические резонансы) \\
			Ось после разрушения & Всегда определима (математическая линия через центр масс) & Перестаёт существовать как физически выделенная линия \\
			\bottomrule
		\end{tabular}
	\end{table}
	
	\subsubsection{Осуществимость и оценка стоимости}
	
	\begin{itemize}
		\item \textbf{Стоимость оборудования:} Криостат (∼\$10 000), вращающийся столик (∼\$20 000), лазерный интерферометр (∼\$15 000), детекторы (∼\$5 000). Итого ∼\$50 000 — в пределах бюджета типичной физической лаборатории или лаборатории материаловедения.
		\item \textbf{Время:} Один месяц на установку, один месяц на измерения, один месяц на анализ.
		\item \textbf{Риск:} Низкий — все компоненты коммерчески доступны, а диски YBCO являются стандартной продукцией.
	\end{itemize}
	
	\subsubsection{Путь к экспериментальной проверке: возможности для сотрудничества}
	
	Предлагаемый эксперимент вполне выполним в существующих лабораториях низких температур и высокоскоростного вращения. Мы предлагаем связаться с исследовательскими группами, которые уже имеют опыт работы с:
	
	\begin{itemize}
		\item \textbf{Вращающиеся криостаты и сверхпроводящие материалы}: например, группа профессора Джона Р. Халла в Вашингтонском университете (ранее в Boeing), или группа по сверхпроводимости в Институте Фрица Габера в Берлине.
		\item \textbf{Высокоточные вращающиеся столики и интерферометрия}: например, группа по гравитационным волнам в MIT (LIGO), имеющая обширный опыт в лазерной интерферометрии и виброизоляции.
		\item \textbf{Высокоскоростная съёмка и исследования фрагментации}: например, группа динамических материалов в Caltech или Институт ударной физики в Университете штата Вашингтон.
		\item \textbf{Российские институты}: Институт физических проблем им. П.Л. Капицы (Москва) имеет давние традиции в физике низких температур и экспериментах с высокоскоростным вращением; Институт физики твёрдого тела (Черноголовка) интенсивно работает с YBCO и другими высокотемпературными сверхпроводниками.
	\end{itemize}
	
	Мы открыты для совместных договорённостей: YUCT может предоставить подробную теоретическую поддержку, включая точные предсказания \(\omega_{\text{crit}}\) для конкретных геометрий дисков и материалов, а также помощь в интерпретации данных. Положительный результат стал бы знаковым открытием: первая демонстрация зависящей от материала геометрии пространства-времени и первое наблюдение координационного фазового перехода во вращающейся системе. Мы приглашаем заинтересованных экспериментаторов связаться с автором по адресу \texttt{alexey@yakushev.eu}.
	
	\subsubsection{Заключение по предлагаемому эксперименту}
	
	Если наблюдаются предсказанные различия между нормальным и сверхпроводящим состояниями, это обеспечит прямое, лабораторное подтверждение утверждения YUCT о том, что геометрия зависит от материала и что ось вращения является эмерджентным свойством координированной материи. Нулевой результат (отсутствие различий) установил бы верхнюю границу для \(K_{\mathrm{eff,mat}} - 1\) для YBCO, ограничивая теорию. Эксперимент прост, доступен и решителен — он превращает парадокс Эренфеста из мысленного эксперимента в эмпирический испытательный стенд для координационной парадигмы.
	
	\subsection{Онтологические следствия: расстояние, время, температура}
	
	Координационная парадигма заставляет нас пересмотреть статус самых фундаментальных физических величин.
	
	\begin{itemize}
		\item \textbf{Расстояние.} В YUCT эффективное расстояние между двумя координационными кластерами \(i\) и \(j\) определяется как \(d_{ij} = c \cdot \langle \tau_{ij} \rangle\), где \(\tau_{ij}\) — среднее время, необходимое для распространения короткого индекса между кластерами. Метрический тензор \(g_{\mu\nu}(x)\) возникает после компактификации 19-мерного координационного поля \(\Psi_{MN}\) и последующего усреднения по внутренним степеням свободы. Следовательно, \textbf{расстояние не является первичной данностью, а является статистической мерой процессов обмена индексами}. Без координированной материи слово «расстояние» не имеет физического смысла.
		
		\item \textbf{Время.} Как показано в Приложении V, собственное время пропорционально увеличению координационной энтропии: \(d\tau \propto dS_{\mathrm{coord}}\). Изолированный агент, не участвующий ни в каком координационном акте, не испытывает времени. \textbf{Время — это мера координационной активности, а не абсолютный внешний параметр}.
		
		\item \textbf{Температура.} Температура является макроскопической проекцией средней координационной ошибки \(\varepsilon\). Универсальный закон ошибок \(\varepsilon = \kappa_c \alpha K_{\mathrm{eff}}^{-\beta}\) вместе с Фундаментальной координационной теоремой (Приложение B) гарантирует, что \(\varepsilon > 0\) для любого конечного \(K_{\mathrm{eff}}\). Следовательно, \textbf{абсолютный нуль температуры недостижим} не из-за технологического ограничения, а потому что идеально координированное состояние невозможно в принципе.
	\end{itemize}
	
	Короче говоря, расстояние, время и температура не являются элементарными строительными блоками реальности. Они являются \textbf{эмерджентными статистическими параметрами} координационной сети, появляющимися только тогда, когда материя координирует свои состояния через обмен индексами и активацию общих словарей.
	
	\section{Возражения и ответы}
	\label{sec:objections}
	
	\paragraph{Возражение 1: «Просто переименование \(\gamma\)»}
	Стандартная ОТО уже использует \(\gamma\) для описания диска. Кажется, что YUCT просто заменяет \(\gamma\) на \(K_{\mathrm{eff}}\) без нового содержания.
	
	\paragraph{Ответ}
	В ОТО \(\gamma\) однозначно определяется \(\omega\) и скоростью света в вакууме. YUCT вводит дополнительный фактор \(K_{\mathrm{eff,mat}}\), который может изменяться независимо (например, с помощью квантовой когерентности). Это приводит к проверяемому отклонению от ОТО: одинаковое \(\omega\) даёт разные длины окружности в зависимости от эффективности координации материала. Более того, отождествление \(K_{\mathrm{eff}} = \gamma\) — это не просто переименование; это конкретная форма, продиктованная универсальными координационными константами \(\Sodd\) и \(\Seven\) и почти равенством \(\Sodd \varphi^2 \approx \pi\). Эти константы обеспечивают более глубокое геометрическое обоснование, отсутствующее в ОТО.
	
	\paragraph{Возражение 2: «Нарушение причинности»}
	Как могут все точки диска знать словарь без сверхсветовых сигналов?
	
	\paragraph{Ответ}
	Словарь предварительно распределён через структуру материала (например, постоянные решётки, спиновое упорядочение). Когда вращение начинается, локальные электромагнитные взаимодействия обеспечивают соблюдение предписанных расстояний. Это не более сверхсветовое, чем твёрдое тело в ньютоновской механике — ограничения встроены, а не передаются в реальном времени.
	
	\paragraph{Возражение 3: «Скрытые переменные, запрещённые теоремой Белла»}
	Предварительно распределённый словарь напоминает скрытую переменную, которую квантовая механика исключает для запутанности.
	
	\paragraph{Ответ}
	Словарь в YUCT не является локальной скрытой переменной в смысле теоремы Белла. Это глобальная, несигнальная корреляционная структура, которая устанавливается офлайн и не может использоваться для передачи информации быстрее света. В квантовой механике само запутанное состояние является такой корреляцией. YUCT просто расширяет это понятие на область классической и релятивистской координации. Теорема Белла запрещает локальные реалистические скрытые переменные, но YUCT не предполагает локального реализма; она предполагает координацию как примитив, и словарь явно нелокален в том смысле, что он предварительно распределён по всей системе. Индекс (само вращение) по-прежнему распространяется причинно. Поэтому YUCT полностью совместима со всеми экспериментальными проверками неравенств Белла.
	
	\paragraph{Возражение 4: «Нет физического механизма для изменения \(K_{\mathrm{eff,mat}}\)»}
	Как можно изменить эффективность координации без изменения \(\omega\)?
	
	\paragraph{Ответ}
	Примеры включают сверхпроводящую когерентность (макроскопическое квантовое состояние), магнитные квантовые осцилляции, индуцированные полем, и оптические решётки. Они напрямую влияют на способность материала поддерживать точные взаимные расстояния, то есть на его эффективность координации. Универсальный показатель масштабирования \(\betaf = 2/3\) из Приложения L обеспечивает количественную связь между объёмом когерентности и ожидаемым изменением \(K_{\mathrm{eff,mat}}\).
	
	\paragraph{Возражение 5: «Формула критической скорости отличается от классической механики разрушения»}
	YUCT предсказывает зависящую от материала критическую скорость, которая может отклоняться от классической формулы. Как это можно согласовать с хорошо установленными инженерными знаниями?
	
	\paragraph{Ответ}
	Для обычных материалов (несверхпроводящих, некогерентных) \(K_{\mathrm{eff,mat}} = 1\) и формула YUCT точно сводится к классическому результату, когда \(\Kcrit\) выражается через предел прочности. Отклонение появляется только при \(K_{\mathrm{eff,mat}} \neq 1\). Таким образом, YUCT не противоречит классической механике; она расширяет её на новые режимы. Классический режим является частным случаем.
	
	\section{Графическое резюме: фазовая диаграмма \(K_{\mathrm{eff}}\)}
	\label{sec:diagram}
	
	\begin{figure}[h]
		\centering
		\begin{tikzpicture}[scale=1.2]
			% Axes
			\draw[->] (0,0) -- (6,0) node[right] {\(K_{\mathrm{eff}}\)};
			\draw[->] (0,0) -- (0,4) node[above] {Геометрия / Физика};
			% Regions
			\fill[blue!10] (0,0) rectangle (2,4);
			\fill[green!15] (2,0) rectangle (4.5,4);
			\fill[red!10] (4.5,0) rectangle (6,4);
			% Labels
			\node at (1,3.5) {Классический};
			\node at (1,3) {\(K_{\mathrm{eff}} = 1\)};
			\node at (1,2.5) {Евклидова};
			\node at (1,2) {Абсолютное время};
			\node at (3.25,3.5) {Релятивистский};
			\node at (3.25,3) {\(K_{\mathrm{eff}} > 1\)};
			\node at (3.25,2.5) {Неевклидова};
			\node at (3.25,2) {Искривлённое пространство};
			\node at (5.25,3.5) {Квантовый};
			\node at (5.25,3) {\(K_{\mathrm{eff}} \to \infty\)};
			\node at (5.25,2.5) {Некоммутативная};
			\node at (5.25,2) {Запутанность};
			% Critical threshold line
			\draw[dashed, thick, orange] (2.5,0) -- (2.5,4);
			\node[orange, rotate=90] at (2.5,2) {\(K_{\mathrm{crit}}\)};
			% Arrows indicating disk evolution
			\draw[thick, ->] (0.5,0.5) to[bend left] (5.5,0.5);
			\node at (3,0.2) {Увеличение \(\omega\) или \(K_{\mathrm{eff,mat}}\)};
			% Vertical line at Keff=1 and Keff=∞ symbolic
			\draw[dashed] (2,0) -- (2,4);
			\draw[dashed] (4.5,0) -- (4.5,4);
			\node at (2, -0.3) {1};
			\node at (4.5, -0.3) {\(\infty\)};
		\end{tikzpicture}
		\caption{Фазовая диаграмма вращающегося диска в YUCT. Один и тот же параметр \(K_{\mathrm{eff}}\) управляет переходом от классического (\(K_{\mathrm{eff}}=1\)) через релятивистский (\(K_{\mathrm{eff}}>1\)) к квантовому (\(K_{\mathrm{eff}}\to\infty\)) поведению. Критический порог \(K_{\mathrm{crit}}\) отмечает границу между устойчивой координацией (слева) и катастрофической потерей координации (справа). Для \(K_{\mathrm{eff,total}} < K_{\mathrm{crit}}\) (за оранжевой линией) диск разрушается.}
		\label{fig:phasediagram}
	\end{figure}
	
	\section{Связь с космическим вращением (Приложение \(\Omega\)) и астрофизическими объектами}
	\label{sec:cosmic}
	
	Те же координационные принципы применимы к вращению Вселенной в целом (см. Приложение \(\Omega\)). Диполь реликтового излучения можно интерпретировать как суперпозицию локального движения и глобального вращения, причём эффективность координации \(K_{\mathrm{eff,cosmic}}\) связана с угловой скоростью \(\omega_{\text{всел}}\). Настоящий анализ вращающегося диска даёт локальный аналог: эффективная метрика для вращающейся Вселенной также описывалась бы зависящей от положения \(K_{\mathrm{eff}}(r)\). Таким образом, проверяемое предсказание зависящей от материала геометрии для вращающегося диска имеет космологический аналог: наблюдаемая амплитуда диполя реликтового излучения может зависеть от крупномасштабной эффективности координации Вселенной. Хотя это и нельзя проверить непосредственно в лаборатории, эта концептуальная связь укрепляет внутреннюю согласованность YUCT.
	
	Нейтронные звёзды (пульсары, магнитары) являются естественными астрофизическими аналогами. Их быстрое вращение (до \(\sim 700\) Гц) и сильные магнитные поля создают экстремальную среду, где эффективность координации определяется гравитацией, ядерными силами и сохранением магнитного потока. Критическое условие разрушения, выведенное в разделе~\ref{sec:critical-disruption}, можно применить для оценки максимальной частоты вращения нейтронной звезды. Используя формулу YUCT с соответствующими параметрами материала (плотность \(\sim 10^{17}\) кг/м³, предел прочности \(\sim 10^{30}\) Па, возможно, \(K_{\mathrm{eff,mat}}\) усиленная сверхтекучестью), получаем максимальную частоту около 1 кГц, что согласуется с наблюдениями. Это даёт независимую перекрёстную проверку рамок.
	
	\subsection{Связь с алгебраическим циклом YUCT: почему геометрия зависит от материала}
	\label{subsec:ehrenfest-algebraic-loop}
	
	Анализ вращающегося диска показал, что эффективная пространственная геометрия — в частности, отношение длины окружности к радиусу — не является абсолютным свойством пространства-времени, а зависит от эффективности координации материала \(K_{\mathrm{eff,mat}}\). Естественно возникает вопрос: является ли \(K_{\mathrm{eff,mat}}\) свободно регулируемым параметром или он ограничен более глубокой алгебраической структурой YUCT?
	
	Ответ кроется в фундаментальных константах, выведенных в Приложении Ferra1.2. Иерархическое разложение координации даёт две универсальные константы:
	\[
	S_{\mathrm{odd}} = \frac{6}{5} = 1.2,\qquad S_{\mathrm{even}} = \frac{4}{5} = 0.8,
	\]
	которые удовлетворяют циклическим соотношениям \(S_{\mathrm{odd}}+S_{\mathrm{even}}=2\) и \(S_{\mathrm{odd}}/S_{\mathrm{even}} = 3/2 = 1/\beta\). Эти константы управляют полуцелым квантованием масс фермионов (\(m_f \propto q^{N_f}\) с \(q = (3/2)^{1/3}\)), а также диктуют предельное поведение координации в макроскопических системах.
	
	Примечательно, что та же алгебраическая структура даёт высокоточное приближение к постоянной окружности \(\pi\). Используя золотое сечение \(\varphi = (1+\sqrt{5})/2\), находим
	\begin{equation}
		S_{\mathrm{odd}} \cdot \varphi^2 = \frac{6}{5} \left( \frac{1+\sqrt{5}}{2} \right)^2 
		= \frac{9+3\sqrt{5}}{5} \approx 3.1416407865\dots
	\end{equation}
	что отличается от \(\pi \approx 3.1415926535\) всего на \(4.8\times10^{-5}\) (относительная ошибка \(1.5\times10^{-5}\)). Это на порядок точнее, чем классическое рациональное приближение \(22/7\).
	
	В YUCT это совпадение не случайно. Эффективное геометрическое отношение \(\pi_{\mathrm{eff}}\) для системы с конечной эффективностью координации \(K_{\mathrm{eff}}\) даётся формулой
	\begin{equation}
		\pi_{\mathrm{eff}}(K_{\mathrm{eff}}) = \pi_{\infty} - \Delta\pi \cdot K_{\mathrm{eff}}^{-\beta},
	\end{equation}
	где \(\pi_{\infty} = \pi\) — предел идеальной координации (\(K_{\mathrm{eff}} \to \infty\)), а \(\Delta\pi\) — константа, связанная с дискретной структурой словаря. Для материала в основном состоянии \(K_{\mathrm{eff}}\) определяется доминирующим иерархическим уровнем, который характеризуется отношением \(S_{\mathrm{odd}}/S_{\mathrm{even}} = 3/2\). Подставляя это в закон масштабирования, получаем \(\pi_{\mathrm{eff}} \approx S_{\mathrm{odd}}\varphi^2\).
	
	\textbf{Последствия для вращающегося диска.} Материальная эффективность координации \(K_{\mathrm{eff,mat}}\), которая модифицирует длину окружности в уравнении (\ref{eq:keff-disk}), таким образом, не является произвольным подгоночным параметром, а коренится в том же алгебраическом цикле, который определяет массы элементарных частиц. Предсказанное отклонение эффективного \(\pi\) от его математического значения при нормальных условиях чрезвычайно мало (поскольку \(K_{\mathrm{eff,mat}}\) велико), что объясняет, почему евклидова геометрия является отличным приближением для повседневного машиностроения. Однако в экстремальных условиях — например, вблизи фазового перехода, где \(K_{\mathrm{eff,mat}}\) резко падает — отклонение может стать обнаружимым. Это даёт прямую связь между парадоксом Эренфеста, проблемой иерархии в физике частиц и предлагаемым лабораторным экспериментом со сверхпроводящим диском (раздел~\ref{sec:experiment}).
	
	Таким образом, зависимость геометрии от материала не является произвольным предположением YUCT, а необходимым следствием той же дискретной алгебраической структуры, которая объединяет спектр фундаментальных фермионов и возникновение непрерывного пространства-времени.
	
	\section{Заключение}
	\label{sec:conclusion}
	
	Мы представили основанную на координации интерпретацию парадокса Эренфеста в рамках YUCT, теперь твёрдо обоснованную универсальными координационными константами \(\Sodd = 1.2\), \(\Seven = 0.8\) и почти равенством \(\Sodd \varphi^2 \approx \pi\). Вращающийся диск рассматривается как система, разделяющая предварительно распределённый словарь взаимных расстояний; вращение действует как индекс, активирующий записи словаря, что приводит к эффективной пространственной метрике, являющейся неевклидовой. Эффективность координации \(K_{\mathrm{eff}}(r)\) непосредственно даёт фактор, на который длина окружности отклоняется от евклидова значения. При \(K_{\mathrm{eff}} \to 1\) геометрия становится евклидовой, а время абсолютным, восстанавливая ньютоновскую физику.
	
	Мы ввели понятия \textbf{координационного центра} (оси вращения) как эмерджентного свойства координированной материи и \textbf{критического разрушения} как катастрофической потери координации, когда \(K_{\mathrm{eff,total}}\) падает ниже зависящего от материала порога \(K_{\mathrm{crit}}\). Мы вывели количественную формулу для критической угловой скорости и показали, что она сводится к классической механике разрушения при \(K_{\mathrm{eff,mat}}=1\). Структура предсказывает, что материалы с повышенной внутренней координацией (например, сверхпроводники) должны демонстрировать другую критическую скорость и изменённую эффективную геометрию — проверяемое отличие как от классической физики, так и от общей теории относительности.
	
	YUCT не противоречит общей теории относительности; вместо этого она даёт более глубокое объяснение происхождения кривизны как эмерджентного свойства высокоэффективной координации, причём иерархические константы обеспечивают количественную основу. Почти равенство \(\Sodd \varphi^2 \approx \pi\) (относительная ошибка \(1.5\times10^{-5}\)) гарантирует, что евклидова геометрия восстанавливается с высокой точностью, когда координация минимальна, а любое отклонение контролируется теми же иерархическими константами. Положительный экспериментальный результат из предлагаемого эксперимента со сверхпроводящим диском не только подтвердил бы YUCT, но и продемонстрировал бы, что геометрией пространства-времени можно управлять через квантовые состояния материала, открывая новый путь для фундаментальной физики.
	
	\begin{thebibliography}{99}
		\bibitem{Ehrenfest1909}
		P. Ehrenfest, ``Gleichf{\"o}rmige Rotation starrer K{\"o}rper und Relativit{\"a}tstheorie,'' \emph{Physikalische Zeitschrift} \textbf{10}, 918 (1909).
		
		\bibitem{Yakushev2026}
		A. V. Yakushev, \emph{Yakushev Unified Coordination Theory (YUCT)}, Zenodo, \url{https://doi.org/10.5281/zenodo.18444598} (2026).
		
		\bibitem{AppendixB}
		Приложение B: Парадокс временной координации -- YUCT.
		
		\bibitem{AppendixL}
		Приложение L: Фрактальное масштабирование ошибок координации -- Универсальный закон от ДНК до космологии.
		
		\bibitem{AppendixR}
		Приложение R: Координационное управление ядерными процессами -- От управляемого распада до холодного синтеза.
		
		\bibitem{AppendixG}
		Приложение G: Объединённая теория координации Якушева -- Фундаментальное решение проблемы квантовой гравитации.
		
		\bibitem{AppendixΩ}
		Приложение \(\Omega\): Глобальное вращение Вселенной и её новый минимальный возраст из радиуса поворота диполя.
		
		\bibitem{AppendixFerra1.2}
		Приложение Ferra1.2: Универсальные координационные константы для упорядоченных и неупорядоченных фаз.
		
		\bibitem{Langevin1935}
		P. Langevin, ``La notion de corps rigide en relativit{\'e},'' \emph{Comptes rendus} \textbf{200}, 1900--1903 (1935).
		
		\bibitem{Gron1995}
		{\O}. Gr{\o}n, ``The relativistic rigid rotating disk,'' \emph{Physics Essays} \textbf{8}, 92--102 (1995).
		
		\bibitem{M87Precession2024}
		Cui, Y. et al., ``Precessing jet in M87: Evidence for a 11-year period from VLBI monitoring 2000–2022,'' \emph{Astrophysical Journal} \textbf{960}, 125 (2024).
		
		\bibitem{NatureAstronomy2025}
		Bright, J. S. et al., ``Transient jet speeds in stellar-mass black holes and neutron stars: The largest sample to date,'' \emph{Nature Astronomy} \textbf{9}, 112–128 (2025).
		
		\bibitem{JetQuenching2006}
		Fender, R. P., Belloni, T. M., \& Gallo, E., ``Towards a unified model for black hole X-ray binaries,'' \emph{Monthly Notices of the Royal Astronomical Society} \textbf{363}, 1105–1118 (2006).
		
		\bibitem{JetQuenching2018}
		Tetarenko, A. J. et al., ``The jet quenching factor in black hole X-ray binaries,'' \emph{Astrophysical Journal} \textbf{862}, 90 (2018).
		
		\bibitem{TDE2016}
		Pasham, D. R. et al., ``A relativistic jet from a tidal disruption event,'' \emph{Astrophysical Journal} \textbf{820}, L18 (2016).
		
		\bibitem{TDE2020}
		Alexander, K. D. et al., ``The tidal disruption event AT2019dsg: A jetted outflow on the rise,'' \emph{Astrophysical Journal Letters} \textbf{894}, L19 (2020).
		
		\bibitem{Yarman2013}
		T. Yarman, ``Radiation from an accelerating neutral body: The case of rotation,'' \emph{European Physical Journal Plus} \textbf{128}, 134 (2013). DOI:10.1140/epjp/i2013-13134-9
		
		\bibitem{Kholmetskii2020}
		A. L. Kholmetskii, T. Yarman, O. Yarman, M. Arik, ``Analyses of Mössbauer experiments in a rotating system: Proper and improper approaches,'' \emph{Annals of Physics} \textbf{418}, 168191 (2020). DOI:10.1016/j.aop.2020.168191
		
		\bibitem{YARK2016}
		A. L. Kholmetskii, T. Yarman, O. Yarman, M. Arik, ``Pound–Rebka result within the framework of YARK theory,'' \emph{Canadian Journal of Physics} \textbf{94}, 806–810 (2016). DOI:10.1139/cjp-2016-0087
		
		\bibitem{Bashkov2000}
		V. Bashkov, M. Malakhaltsev, ``Geometry of rotating disk and the Sagnac effect,'' arXiv:gr-qc/0011061 (2000).
		
	\end{thebibliography}
	
\end{document}